\documentclass[11pt,a4paper]{article}

% ============================================================
% Uniform MTDF preamble -- identical across all papers
% ============================================================
\usepackage[utf8]{inputenc}
\usepackage[T1]{fontenc}
\usepackage{lmodern}
\usepackage[english]{babel}
\usepackage[a4paper,margin=2.2cm]{geometry}
\usepackage{amsmath,amssymb,amsthm}
\usepackage{graphicx}
\usepackage{float}
\usepackage{booktabs}
\usepackage{array}
\usepackage{longtable}
\usepackage{tabularx}
\usepackage{xltabular}
\usepackage{multirow}
\usepackage{caption}
\usepackage[dvipsnames,table]{xcolor}
\usepackage{mdframed}
\usepackage{parskip}
\usepackage{ragged2e}
\usepackage[hidelinks,breaklinks=true]{hyperref}
\usepackage{microtype}
\usepackage{url}
% Allow \path{} to break more aggressively inside narrow table columns:
\makeatletter
\g@addto@macro\UrlBreaks{\do\_\do\-\do\.\do\/\do\:}
\makeatother

% Colours matching the HTML theme
\definecolor{accent}{HTML}{1A4B8A}
\definecolor{callout-bg}{HTML}{FBFBFB}
\definecolor{tablehead}{HTML}{F0F0F0}

% Prevent any table from overflowing
\setlength{\tabcolsep}{4pt}
\renewcommand{\arraystretch}{1.15}

% Caption style (compact, italic, small like the HTML)
\captionsetup{font=small,labelfont=bf,labelsep=period,justification=centerfirst}

% Hyperref metadata placeholders (filled per-paper below)
\hypersetup{
  pdftitle={Mesche's Tensor Dynamics Framework (MTDF): A Dynamic Field Theory of Cosmic Structure},
  pdfauthor={Ingo Mesche},
  pdfsubject={MTDF - Mesche Time-Dilation Field Theory},
  colorlinks=false
}

% Prevent overfull hbox from long URLs/DOIs in bibliography
\sloppy
\emergencystretch=3em

% Tolerant line breaking so long identifiers (Python filenames etc.) don't
% produce large overfull hboxes in narrow table columns. These settings keep
% all content on-page while slightly relaxing right-margin strictness.
\tolerance=1200
\hbadness=2000
\hyphenpenalty=150
\exhyphenpenalty=50

% ============================================================
\title{Mesche's Tensor Dynamics Framework (MTDF): A Dynamic Field Theory of Cosmic Structure}
\author{Ingo Mesche\\ \small Independent Researcher, Malta \\ \small \texttt{imesche@outlook.com}}
\date{V75 / July 2026}

\begin{document}
\maketitle
\begin{center}\itshape Version 75 (Master Paper). Corrected relativistic architecture (frozen CON-4B unified action), corrected covariant growth sector, and sealed v2 full-likelihood results. Produced non-destructively from V74 under MTDF\_\allowbreak{}REWRITE\_\allowbreak{}MANIFEST\_\allowbreak{}V3.\end{center}




\begin{abstract}
Mesche's Tensor Dynamics Framework (MTDF) is a field based alternative to the latent sector interpretation of modern cosmology. It aims to reproduce the background expansion history and a wide range of dynamical phenomena without invoking cold dark matter or a cosmological constant as new substances. This master paper presents the theory as a coherent physical proposal, states the minimal parameter set and its provenance, and formalises a strict evaluation protocol that separates calibration anchors from independent validation targets. It then summarises current V18 validation outcomes on the vector likelihood set (Pantheon+ supernovae, DESI Y1 BAO, cosmic chronometer H(z), DR16 f$\sigma$$_{8}$ growth) and reports the CMB distance prior strictly as a diagnostic. A sharp, environment-dependent transition in SNe residuals is identified at z = 0.04 $\pm$ 0.005 across three independent void catalogues, with $\Delta$$\chi$$^{2}$ = 36--39 relative to a single-regime fit (nominal p < 0.001, diagonal-covariance merged sample; see Companion Part I for covariance caveats and replication requirements). This scale (\~{}170 Mpc comoving) coincides with the measured southern extent of the local underdensity (B\"ohringer et al. 2020); the strain-lever length scale $\beta$ = 22.7 Mpc is a distinct, smaller field scale, and any quantitative connection between the two is a structural conjecture of the framework discussed in Section 10.1. Fitting the full Planck 2018 TTTEEE+lensing likelihood with the corrected class\_\allowbreak{}mtdf solver (v2 chains; corrected covariant $\mu$(a) in newtonian gauge, with the high-Q carrier represented by its CDM-like limit, a proxy whose scope is stated in Section 7.3) yields a broad posterior on the MTDF coupling, k\textsubscript{f} = 0.685 $\pm$ 0.500, in which both k\textsubscript{f} = 0 ($\Lambda$CDM) and k\textsubscript{f} = 1 (full coupling) lie within about 1.4$\sigma$ of the mean: the full Planck likelihood neither favours nor excludes the coupling, and this dataset establishes consistency, not support. Standard-parameter shifts are negligible and the CMB-side H\textsubscript{0} is unchanged. Under the same constraints, the corrected covariant growth sector enhances late-time growth by approximately 2.5 percent at fixed parameters; the v2 chains give $\sigma$\textsubscript{8} = 0.8265 $\pm$ 0.0062, about 1.9$\sigma$ (posterior units) above the same-pipeline $\Lambda$CDM baseline of 0.8101, so the corrected theory mildly worsens, rather than eases, the comparison with published weak-lensing clustering amplitudes; the residual weak-lensing tension is shared with $\Lambda$CDM, and the only pre-registered route by which this posture could be revised is the forward-modelling programme (FWD-S8). The S\textsubscript{8} = 0.802 relief reported in V74 was a gauge artifact and is formally retired. Against the completed matched $\Lambda$CDM control chain, the bounded best fits give $\Delta$$\chi$$^{2}$ (MTDF $-$ $\Lambda$CDM) = $-$1.29: the full Planck likelihood does not distinguish the corrected MTDF from $\Lambda$CDM. This is a maximum-likelihood comparison, not a Bayes factor or Bayesian evidence calculation. The paper emphasises falsifiability: it enumerates near term tests that can confirm or refute MTDF beyond the current dashboard scope. The framework's known weaknesses, including the post-hoc origin of the environmental equation-of-state coupling and the fact that the stress-matter coupling $\alpha$ is currently motivated by void dynamics rather than independently measured, are stated explicitly and paired with falsification conditions. A separate gravity-sector companion paper reports quantitative lensing and local dynamics tests (galaxy-galaxy lensing, strong lens time delays, and Solar System screening) using the same frozen parameter set; see \textit{MTDF\_\allowbreak{}03 (Companion Part II)}~[see companion paper].
\end{abstract}

\tableofcontents

\begin{mdframed}[linecolor=accent,linewidth=0.5pt,backgroundcolor=callout-bg,leftline=true,rightline=true,topline=true,bottomline=true,innerleftmargin=8pt,innerrightmargin=8pt,innertopmargin=6pt,innerbottommargin=6pt]
\textbf{Claim status.}
This paper distinguishes between:
(i) calibration anchors,
(ii) validation targets,
(iii) diagnostic consistency checks, and
(iv) speculative extensions.
Only results explicitly labelled as validation targets are used as evidence for the core MTDF claim. Diagnostic and speculative sections are included to define future tests, not to claim established confirmation.

\textbf{Claim classes (frozen ledger).} Every claim of the corrected architecture carries exactly one of eight classes from the frozen claim ledger: VALIDATED (derived and numerically confirmed against data or an independent implementation), DERIVED (theorem or derivation grade, not yet numerically confronted), DECLARED (postulate in honest form, with forcing ancestry stated), CONDITIONAL (named calculation owed), PREDICTION-REGISTERED (dated, falsifiable, awaiting data), CONTESTED (prediction stated, data unsettled), OPEN (identity unknown), and RETIRED (tombstoned, with a pointer to the killing record). Where this paper asserts a class, it is the class recorded in the claim-status scheme above.
\end{mdframed}

\section{Motivation and design principles}
\label{sec:sec1}
Contemporary cosmology describes most observed gravitational phenomena by introducing a dominant but unseen sector. In the standard model of cosmology, these components are parameterised as cold dark matter and dark energy. MTDF proposes an alternative interpretation: the phenomena attributed to the latent sector emerge from the dynamics of a physically defined spacetime field that carries stress energy and couples to baryonic matter through a single coupling parameter.

\begin{mdframed}[linecolor=accent,linewidth=0.5pt,backgroundcolor=callout-bg,leftline=true,rightline=true,topline=true,bottomline=true,innerleftmargin=8pt,innerrightmargin=8pt,innertopmargin=6pt,innerbottommargin=6pt]
\textbf{Design constraints for a high standard theory proposal.} MTDF is formulated under four constraints: (i) minimal free parameters; (ii) explicit parameter provenance; (iii) strict separation of calibration anchors from validation targets; (iv) falsifiability with clear failure modes. The V18 era artefacts operationalise these constraints for empirical evaluation.
\end{mdframed}

It is essential to distinguish between two different claims: (a) \emph{data fit} within a chosen likelihood set, and (b) \emph{physical explanation} that remains coherent across scales. The dashboard is primarily about (a). This master paper addresses (b), while importing the dashboard protocol to ensure that any claimed fit is not achieved by circular calibration.

\section{Core objects and governing assumptions}
\label{sec:sec2}
MTDF treats the gravitational sector as augmented by a field degree of freedom whose stress energy contributes to dynamics. The field is introduced to remain compatible with classical limits while providing a mechanism for scale dependent behaviour. In the non relativistic regime, the field acts as an additional contribution to the effective potential that influences orbital velocities and structure formation, while in the cosmological regime it modifies the effective expansion history through its energy density and stress.

\subsection{Minimal parameter philosophy}
The framework is constrained to a small set of fundamental parameters that are either measured from independent observational anchors or established from theoretically motivated universal constants. The operational intent is that most quantities used in predictions are derived, not tuned.

\subsection{Calibration anchors versus validation targets}
In the V18 strict protocol, some data are permitted to set fundamental scales (anchors). These are excluded from validation scoring. All other reported goodness of fit values are computed strictly on targets that are treated as independent of the calibration step. This is the key correction relative to earlier single layer evaluations.

\subsection{Mathematical objects, stress definition, and energy}
This section expands the core mathematical objects that appear throughout MTDF. The intent is to make the field definitions explicit, while keeping the strict V18 protocol separation between calibration anchors and validation targets.

\begin{mdframed}[linecolor=accent,linewidth=0.5pt,backgroundcolor=callout-bg,leftline=true,rightline=true,topline=true,bottomline=true,innerleftmargin=8pt,innerrightmargin=8pt,innertopmargin=6pt,innerbottommargin=6pt]
\subsubsection{Core idea in one sentence}
We propose that the universe is permeated by dynamic spacetime characterized by the dimensionless strain tensor S\textsubscript{$\mu$$\nu$} that:

\begin{itemize}
  \item \textbf{Stores energy} through elastic deformation with energy density u\textsubscript{tensor} = (E/\allowbreak{}2) Tr(S\textsubscript{$\mu$$\nu$} S\textsubscript{$\mu$$\nu$})
  \item \textbf{Exhibits stress tensor gradients} that wrap around massive objects and extend throughout cosmic space
  \item \textbf{Undergoes stress accumulation and relaxation cycles} over characteristic time $\tau$ = 13.0 Gyr
  \item \textbf{Redistributes energy} rather than destroying it, maintaining conservation laws
  \item \textbf{Exhibits bipolar behaviour} with attractive and repulsive components at different scales
\end{itemize}

\begin{mdframed}[linecolor=accent,linewidth=0.5pt,backgroundcolor=callout-bg,leftline=true,rightline=true,topline=true,bottomline=true,innerleftmargin=8pt,innerrightmargin=8pt,innertopmargin=6pt,innerbottommargin=6pt]
\subsubsection{Core MTDF Parameters (provenance classified in Section 5.0)}
\textbf{Reference:} Complete parameter values and derivations are summarised in Section 5 and in the validation workbook accompanying the release.
\end{mdframed}
\end{mdframed}

MTDF represents the additional degree of freedom through a dimensionless strain field. In compact notation one may write:

\[
\tilde{S}_{\mu\nu} = \beta\,\nabla_\mu \xi_\nu
\]

Stress is defined through a constitutive relation with an effective elastic modulus:

\[
\Sigma_{\mu\nu} = E\,\tilde{S}_{\mu\nu}
\]

\begin{mdframed}[linecolor=accent,linewidth=0.5pt,backgroundcolor=callout-bg,leftline=true,rightline=true,topline=true,bottomline=true,innerleftmargin=8pt,innerrightmargin=8pt,innertopmargin=6pt,innerbottommargin=6pt]
\subsubsection{Construction versus postulate}
The primary field degree of freedom in MTDF is the displacement-like field $\xi$\textsubscript{$\nu$}. The strain tensor S\textsubscript{$\mu$$\nu$} is defined from its covariant gradient, and the elastic stress tensor $\Sigma$\textsubscript{$\mu$$\nu$} follows from the constitutive relation $\Sigma$\textsubscript{$\mu$$\nu$} = E S\textsubscript{$\mu$$\nu$}. Thus $\Sigma$\textsubscript{$\mu$$\nu$} is not introduced as an independent field-strength tensor or as an additional matter component. It is the stress response of the spacetime medium associated with the strain field.
\end{mdframed}

The associated field energy density and an effective negative pressure contribution can be written schematically as:

\[
u_\mathrm{tensor} = \frac{E}{2}\,\tilde{S}_{\mu\nu}\,\tilde{S}^{\mu\nu},\qquad
p_\mathrm{tensor} \approx -\frac{1}{3}u_\mathrm{tensor}
\]

\subsection{Evolution equation and relaxation time}
In the simplest non relativistic implementation used for many intuition building arguments, MTDF employs a damped wave type evolution for an effective potential-like field variable. A representative form is:

\[
\frac{\partial^2 \phi}{\partial t^2} - c_s^2\,\nabla^2\phi + \gamma\,\frac{\partial \phi}{\partial t} = 8\pi G\rho
\]

Here \(\rho\) denotes the baryonic source term, \(c_s\) is an effective propagation speed for stress perturbations, and \(\gamma\) sets the relaxation rate. In many analytic arguments \(\gamma\) is treated as the inverse of a cosmic scale relaxation time \(\tau\), which is part of the minimal parameter set.

\begin{mdframed}[linecolor=accent,linewidth=0.5pt,backgroundcolor=callout-bg,leftline=true,rightline=true,topline=true,bottomline=true,innerleftmargin=8pt,innerrightmargin=8pt,innertopmargin=6pt,innerbottommargin=6pt]
\subsubsection{On analogies and what is not claimed}
The stress-strain behavior of spacetime ($\Sigma$\textsubscript{$\mu$$\nu$} = E S\textsubscript{$\mu$$\nu$}) manifests dynamical patterns with direct analogues in condensed matter systems.

Such parallels provide physical motivation for the stress-tensor analogy used in MTDF, illustrating how material physics principles scale to cosmic phenomena when applied to spacetime itself with elastic modulus E = 9.1$\times$10$^{-}$$^{1}$$^{0}$ Pa.
\end{mdframed}

\subsection{Regime transition and critical acceleration}
The framework is constructed to recover the classical limit at high accelerations while producing measurable departures in the low acceleration regime. A commonly used summary quantity is a critical acceleration scale that depends on the fundamental length scale \(\beta\):

\[
a_\mathrm{crit} \sim \alpha\,\frac{c^2}{\beta}
\]

This scale is used as a compact way to discuss why deviations are negligible in Solar System tests yet become relevant in galaxies and the late time large scale environment. Precise operational definitions depend on the specific observable being computed and are therefore treated as derived quantities in the V18 artefacts.

\subsection{Local rigidity and scale hierarchy}
The acceleration-dependent relaxation (Section 2.4) provides one mechanism for suppression in high-acceleration environments. A second, independent mechanism arises from the substrate's wavelength-dependent stiffness. The effective stiffness of the MTDF substrate at wavelength \(\lambda\) is:

\[
K_{\mathrm{eff}} = E\,\beta^2\,(2\pi/\lambda)^2
\]

At laboratory scales (\(\lambda \sim 0.1\) m), \(K_{\mathrm{eff}} \sim 10^{40}\) Pa. At solar system scales (\(\lambda \sim 1\) AU), \(K_{\mathrm{eff}} \sim 10^{14}\) Pa. These vastly exceed any stress that laboratory or solar-system sources can produce.

This rigidity is independent of damping: even in a zero-acceleration environment, the substrate resists short-wavelength perturbations because the strain-lever length scale \(\beta = 22.7\) Mpc amplifies strain relative to displacement (\(\tilde{S} = \beta\,\nabla\xi\)). The same lever that makes MTDF responsive at galactic scales (\(\beta k \ll 1\)) makes it rigid at laboratory scales (\(\beta k \gg 1\)).

MTDF therefore exhibits a natural scale hierarchy controlled by the dimensionless product \(\beta k\), where \(k = 2\pi/\lambda\):

\begin{table}[H]
\centering
\footnotesize
\begin{tabularx}{\linewidth}{l>{\RaggedRight\arraybackslash}X>{\RaggedRight\arraybackslash}X>{\RaggedRight\arraybackslash}X>{\RaggedRight\arraybackslash}X}
\toprule
\textbf{Regime} & \textbf{\(\beta k\)} & \textbf{Stiffness} & \textbf{Damping} & \textbf{Outcome} \tabularnewline
Laboratory & \(\sim 10^{24}\) & Rigid & Damped & Undetectable \tabularnewline
Solar system & \(\sim 10^{20}\) & Rigid & Damped & Undetectable \tabularnewline
Galaxy & \(\sim 10^{-5}\) & Soft & Undamped & Active \tabularnewline
Cosmological & \(\sim 10\) & Soft & Undamped & Active \tabularnewline
\bottomrule
\end{tabularx}
\end{table}

This hierarchy follows directly from the fundamental parameters and requires no tuning. The dual mechanism (damping + rigidity) provides a complete explanation for the absence of MTDF effects in all local experiments, including precision electromagnetic measurements with superconducting cavities (\(Q > 2 \times 10^{11}\)), without requiring any additional assumptions.

\section*{2A. The relativistic completion: the frozen unified action}
\addcontentsline{toc}{section}{2A. The relativistic completion: the frozen unified action}
\label{sec:sec2a}
This section states the canonical relativistic architecture of MTDF, frozen as an immutable snapshot on 2026-07-09 at the end of a construction programme; the resulting claim classes and safety table are stated in full in this paper and in Companion Part II. Everything in this section is the source of truth for the rest of the suite: where an earlier effective description in this paper or its companions differs from the frozen action, the frozen action governs, and the earlier description is retained only as a labelled effective or historical parameterisation.

\subsection*{2A.1 The action}
\addcontentsline{toc}{subsection}{2A.1 The action}
The frozen MTDF action, designated CON-4B in the construction programme, is

\[
S \;=\; S_{\mathrm{EH}} \;+\; \int d^4x\,\sqrt{-g}\;\Big[\, \mathcal{L}_{\mathrm{core}} \;+\; \mathcal{L}_{J}(\phi) \;+\; \mathcal{L}_{A}(A^\mu) \,\Big] \;+\; S_{\mathrm{m}}\big[\tilde{g}_{\mu\nu},\,\Psi\big],
\]

with the following ingredients.

\begin{itemize}
  \item \textbf{Elastic solid core} \(\mathcal{L}_{\mathrm{core}} = -U(h^{AB})\): a Carter-Quintana elastic solid whose material fields supply the background current \(u_{\mathrm{bg}}\) (the dark-energy role, \(w \approx -1\)). The strain is built from the displacement gradient, \(\tilde{S} = \beta\,\nabla\xi\); because \(U\) depends on the material fields only through the shift-invariant induced metric \(h^{AB}\), the phonons are massless Goldstone modes of broken translations. This exclusion of any displacement mass is load-bearing for the growth sector (Section 2B).
  \item \textbf{One unified scalar} \(\phi\) with a shift-symmetric constitutive function \(J(Q)\), \(Q = |\nabla\phi|^2/a_{\mathrm{MTDF}}^2\). Its low-\(Q\) limit is the declared aquadratic (A0 cubic) response \(J \to (2/3)Q^{3/2}\), which yields the deep-MOND relation \(g = \sqrt{g_{\mathrm{bar}}\,a_{\mathrm{MTDF}}}\) and the baryonic Tully-Fisher scaling \(V_{\mathrm{flat}}^4 = G M_b\, a_{\mathrm{MTDF}}\) (\(V \sim M^{1/4}\)); the radial acceleration relation is validated against SPARC at \~{}0.11 dex scatter. Its high-\(Q\) limit is linear, a CDM-like cosmological carrier. The acceleration scale is the solid-anchored \(a_{\mathrm{MTDF}} = \kappa_a \sqrt{G\,u_{\mathrm{bg}}}\), not a free parameter: the anchor relation is derived, while the O(1) constant \(\kappa_a\) remains a named open research item.
  \item \textbf{A unit-timelike vector} \(A^\mu\) (Einstein-aether class kinetic term) with the coefficient condition \(c_{13} = c_1 + c_3 = 0\), which aligns the gravitational-wave cone with the photon cone exactly (\(c_T = 1\), validated against GW170817).
  \item \textbf{Matter and photons on the disformal metric} \(\tilde{g}_{\mu\nu} = A^2(Q)\, g_{\mu\nu} + B(Q)\, A_\mu A_\nu\), with \(B = -4\,S_{A0}\) and \(A^2 = e^{B/2}\), where \(S_{A0}\) is the A0/\allowbreak{}MOND scalar potential. This algebraic disformal coupling delivers lensing = dynamics (\(\eta = 1\)) without a new propagating mode.
\end{itemize}

\textbf{Degree-of-freedom count (corrected attribution).} The propagating content is \textbf{9 degrees of freedom}: graviton 2 + solid phonons 3 + unified scalar 1 + vector 3. The extra degree of freedom relative to a naive count is the vector, not a second scalar; the unification of the A0 postulate with the carrier removed a conceptual role-redundancy, not a propagating mode.

\subsection*{2A.2 Division of labour}
\addcontentsline{toc}{subsection}{2A.2 Division of labour}

\begin{table}[H]
\centering
\footnotesize
\begin{tabularx}{\linewidth}{>{\RaggedRight\arraybackslash}X>{\RaggedRight\arraybackslash}X}
\toprule
\textbf{Observable} & \textbf{Owner} \tabularnewline
\midrule
Background / \(u_{\mathrm{bg}}\) & solid core \tabularnewline
\(a_{\mathrm{MTDF}}\) anchor & solid \(u_{\mathrm{bg}}\) (\(a_{\mathrm{MTDF}} = \kappa_a\sqrt{G u_{\mathrm{bg}}}\)) \tabularnewline
Galactic dynamics (RAR/\allowbreak{}BTFR) & unified scalar, low-\(Q\) branch (= A0 cubic) \tabularnewline
Lensing (galaxy + cluster) & unified scalar low-\(Q\) on \(\tilde{g}\) (\(\eta=1\)) + vector (GW-safe disformal) \tabularnewline
Linear growth & corrected covariant \(\mu(a)\) (Section 2B) \tabularnewline
CMB peaks & unified scalar high-\(Q\) (CDM-like carrier) + baryons \tabularnewline
Cluster abundance & linear amplitude (carrier + \(\mu\)); unified-scalar nonlinear boost, tau-guard-moderated \tabularnewline
GW speed & vector kinetic (\(c_{13}=0 \Rightarrow c_{\mathrm{GW}}=c_{\mathrm{EM}}\)) \tabularnewline
PPN / local & low-\(Q \to\) GR at high acceleration; disformal \(B\to 0\); solid rigid at \(\beta k \gg 1\) \tabularnewline
H\textsubscript{0} stratification & EFE (\(k_f\)) + environment (extension) \tabularnewline
\bottomrule
\end{tabularx}
\end{table}

Division-of-labour table of the frozen action: one owner per observable, no empty row, no double-counting.

\textbf{Scope restriction (declared).} The tau-guard, the relaxation-clock moderation \(B(t_{\mathrm{coh}}/\tau)\) of the nonlinear response, operates at \emph{cluster scale only}. This is a DECLARED scope restriction, an authorial choice with stated ancestry (a structural analysis, an empirical cross-check, and an explicit authorial decision recorded in the construction log), not a theorem falling out of the action. At galaxy scale the response is unmodulated: the radial acceleration relation is universal, with no assembly dependence, as validated across SPARC. The galaxy-scale assembly test survives as a conditional prediction with a named revival dataset (WALLABY-era young dwarfs with resolved rotation curves and star-formation histories).

\subsection*{2A.3 Health statement}
\addcontentsline{toc}{subsection}{2A.3 Health statement}
The architecture is \textbf{clean at the classical and quadratic level}, within the derived \(K_B\) window and under the stated imported Einstein-aether/\allowbreak{}AeST health assumptions: the cone-coincidence condition is inherited from the AeST construction, and the vector-sector bounds are inherited from the published Einstein-aether literature (Oost, Mukohyama \& Wang 2018; Yagi et al. 2014; Foster \& Jacobson 2006). \(K_B\) is upper-windowed only; a possible strong-coupling lower floor for small \(K_B\) is a named post-paper research item. This wording is deliberate: the health statement claims exactly what was checked, no more.

\textbf{The internal branch switch introduces no instability and no preferred wavenumber.} The unified \(J\)-function passes from the low-\(Q\) A0 cubic to the high-\(Q\) carrier through an elliptic, non-oscillatory region: for any monotone MOND interpolation both quadratic-action eigenvalues remain positive (\(\mu(Q) > 0\), \(\mu + 2Q\mu' > 0\)) and the operator-anisotropy ratio stays in \((1/2, 1]\), so there is no ghost, no gradient instability, and no resonant \(k_*\) anywhere in the frozen action. The transition width (corrections of \~{}5--24 percent one decade either side of \(Q_* \sim 1\)) is empirically pinned by the validated RAR shape, not a free residue. A published AeST oscillatory regime arose in that construction's \(\tilde{\mu}\)-function; the frozen \(J\) avoids it by the monotone-MOND structure the RAR pins, and where MTDF's function differs from that construction is exactly where theirs oscillated and ours cannot.

\subsection*{2A.4 What is imported, and why the import is minimal}
\addcontentsline{toc}{subsection}{2A.4 What is imported, and why the import is minimal}
MTDF imports from the AeST class exactly one structural element it could not supply natively: the vector \(A^\mu\). The necessity is derived, not chosen: a kinetic-free composite disformal coupling would set \(c_{\mathrm{photon}}^2 = 1 - B/A^2\), splitting the photon and GW cones by \~{}9 $\times$ 10\textsuperscript{$-$7} and violating the GW170817 bound by roughly nine orders of magnitude. The vector kinetic term is precisely the degree of freedom that removes this, and its one irreplaceable job (ghost-free cone coincidence) is a ratio condition, \(c_{13} = 0\), independent of the vector's overall scale. The four independent fences on the vector kinetic normalisation (Einstein-aether pulsar band; vector-sector health; the carrier role with BBN/\allowbreak{}PPN bounds; the \(c_T = 1\) surface) force \(K_B\) into the window \((0, \sim 2.5\times10^{-5}]\), the near-GR-aether limit. MTDF is comfortable exactly there because the scalar carries all the load: the vector is an alignment field, not a workhorse. \(K_B\) is stated as windowed, never as an unconstrained O(1) constant.

\textbf{The unification decision paid twice.} The single-scalar unification (one \(J\)-function owning both the galactic A0 response and the cosmological carrier) was taken for economy. It then proved load-bearing for safety: had the carrier or the lensing role demanded an O(1) vector kinetic term while BBN and PPN demanded \~{}10\textsuperscript{$-$5}, the window above would have closed to an empty set and the architecture would have failed post-freeze. What kept the window open is that the scalar carries everything, so no role needs a large vector. Decisions taken for economy that later carry a safety proof are recorded as such in the construction ledger.

\subsection*{2A.5 Open research frontier}
\addcontentsline{toc}{subsection}{2A.5 Open research frontier}
\begin{itemize}
  \item \(\kappa_a\): the O(1) anchor constant in \(a_{\mathrm{MTDF}} = \kappa_a\sqrt{G u_{\mathrm{bg}}}\).
  \item \(\beta_{\mathrm{eos}}\): first-principles origin of the frozen value 0.573.
  \item \(K_B\): its placement inside the derived window \((0, \sim 2.5\times10^{-5}]\), plus the possible strong-coupling lower floor for small \(K_B\).
  \item Full-action Boltzmann implementation of the 9-degree-of-freedom perturbation system (unified scalar + vector perturbations); the present cosmological validation uses the corrected effective implementation with the carrier's CDM-like limit as proxy (scope statement in Section 7.3).
  \item A single constitutive calculation may plausibly deliver \(\kappa_a\), \(\beta_{\mathrm{eos}}\), and \(K_B\) jointly; this is scoped as post-paper work.
\end{itemize}

\begin{mdframed}[linecolor=accent,linewidth=0.5pt,backgroundcolor=callout-bg,leftline=true,rightline=true,topline=true,bottomline=true,innerleftmargin=8pt,innerrightmargin=8pt,innertopmargin=6pt,innerbottommargin=6pt]
\subsubsection*{Relation to the effective descriptions elsewhere in this paper}
The enhancement law \(v_c^2(r) = (GM/r)\,[1 + \alpha/(1 + r/\beta)]\) used in Section 4.1 and in the validation pipeline is the historical effective parameterisation of the galactic response. The canonical mechanism is the low-\(Q\) branch of the frozen unified action, whose aquadratic response produces the curved, slope-1/\allowbreak{}2 low-acceleration bend of the radial acceleration relation that the constant-enhancement limit could not: the shape-level question left open in earlier versions of this paper is resolved by the unified action, at the price stated honestly in the claim ledger, namely that the low-\(Q\) aquadratic form is a declared postulate (honest form, with the RAR/\allowbreak{}BTFR as its forcing ancestry), not a derivation from the elastic core.
\end{mdframed}

\section*{2B. Canonical growth sector: the corrected covariant \(\mu(a)\)}
\addcontentsline{toc}{section}{2B. Canonical growth sector: the corrected covariant $\mu$(a)}
\label{sec:sec2b}
This section is the canonical record of the MTDF late-time growth mechanism. It supersedes all earlier growth-sector descriptions in this suite.

\subsection*{2B.1 Background equation and field identity}
\addcontentsline{toc}{subsection}{2B.1 Background equation and field identity}
The isotropic background stress amplitude obeys

\[
\ddot{S}_0 \;+\; 3H\,\dot{S}_0 \;+\; m^2 S_0 \;=\; \xi\,\rho_m .
\]

Here \(S_0\) denotes the isotropic background amplitude of the strain stress, \(\tilde{S} = \beta\,\nabla\xi\), a shift-invariant (gradient-built) quantity, and not the Goldstone displacement \(\xi\). The coefficient \(m^2\) is the effective restoring curvature of the shift-invariant strain-energy \(U(h^{AB})\) (the equation-of-state / residual-tension bulk modulus); it governs the temporal relaxation of the background amplitude and renormalises the phonon gradient (sound-speed) energy at order \(k^2\). It is not a mass for the displacement \(\xi\): the material shift symmetry keeps the phonons gapless (\(\omega = c_s k\)), and the Goldstone-forbidden displacement mass \(m = 1/\beta\) is absent from \(U(h^{AB})\). The late-time modification \(\mu(a) - 1 = \lambda_{\mathrm{MTDF}}\, T(a/a_t)\), with

\[
\lambda_{\mathrm{MTDF}} \;=\; \frac{(1-\beta_{\mathrm{eos}})^2}{1+\alpha} \;=\; 0.0793,
\]

is a background (\(k\)-independent) effective-gravity effect; its amplitude is set by \(\lambda_{\mathrm{MTDF}}\) and does not depend on \(m^2\), which fixes only the timing of \(S_0\)'s relaxation. This \(m^2\) is also distinct from the retired wavenumber-dependent mediator mass of the abandoned Helmholtz-filtered growth branch: no mass term for the displacement, and no filter scale, appears anywhere in the frozen action.

\subsection*{2B.2 One coupling, two observables}
\addcontentsline{toc}{subsection}{2B.2 One coupling, two observables}
The same derived coupling \(\lambda_{\mathrm{MTDF}}\) sets both the late-time growth modification and the early field-energy injection: \(\mu(a) = 1 + \lambda_{\mathrm{MTDF}}\,T(a/a_t)\) with \(T(x) = x^\alpha/(1+x^\alpha)\), and \(f_{\mathrm{kick}} = \lambda_{\mathrm{MTDF}}/24 = 0.0033\). One coupling therefore unifies the growth sector and the EFE; neither is tuned per-observable. Numerically \(\mu > 1\) at all redshifts (1.026 at z = 2, 1.036 at z = 1, 1.053 at z = 0), and the integrated effect on the growth history is an enhancement of approximately +1.0 percent at z = 2 rising to +2.5 percent at z = 0, at fixed parameters. The mechanism is CMB-inert: \(\mu(z_{\mathrm{rec}}) = 1.00002\).

\subsection*{2B.3 Direction of the effect, and mechanism separation}
\addcontentsline{toc}{subsection}{2B.3 Direction of the effect, and mechanism separation}
\textbf{The corrected covariant \(\mu(a)\) enhances late-time linear growth.} The growth-suppression direction reported in earlier versions of this suite was a gauge artifact (a modification applied to the synchronous-gauge 00-source broke a near-cancellation and flipped the sign; the corrected implementation applies the covariant \(\mu\) in newtonian gauge) and is formally retired, together with the $\sigma$\textsubscript{8} = 0.790 result it produced.

\textbf{Two distinct mechanisms, never conflated.} (a) \emph{Late linear growth:} the corrected \(\mu(a)\), effectively active below the transition near \(z_t = 0.74\), produces the \~{}+2.5 percent enhancement and the clustering-amplitude pressure discussed in Section 7.3. It cannot cause structure at z \~{} 10 and no statement in this suite attributes early massive galaxies to it. (b) \emph{Early nonlinear collapse:} the non-analytic response of the declared A0 cubic is the mechanism that points in the direction of early structure (the JWST-direction effect). It is the same declared postulate that gives the RAR and BTFR, not a fitted knob; cluster-scale claims built on it always carry the cluster-only tau-guard named as the guard. A registered growth-history signature (GH-1) freezes the predicted trajectories: the f$\sigma$\textsubscript{8} excess is already \~{}2 percent at z = 1.5 and becomes more pronounced below the \(\mu(a)\) transition near \(z_t = 0.74\); it does not begin at the transition.

\section{Cosmological background dynamics}
\label{sec:sec3}
At background level, MTDF is evaluated by comparing its predicted distance and expansion relations to vector likelihood datasets. The present implementation in the V18 artefacts uses a dedicated pipeline to compute the background observables and their residuals against each dataset.

\subsection{Dual epoch mechanism for the Hubble tension}
\emph{Notation.} The symbol $\beta$\textsubscript{eos} denotes the dimensionless equation-of-state parameter (0.573); it is distinct from the physical strain-lever length scale $\beta$ = 22.7 Mpc introduced in Section 2. Both are listed separately in Table 1.

The theory separates the Hubble tension into an early epoch contribution and a late epoch contribution. The early epoch mechanism is encoded in a small derived energy injection fraction, defined from the fundamental parameter pair ($\alpha$, $\beta$\textsubscript{eos}), and is used to assess consistency against the CMB distance prior as a diagnostic. The late epoch mechanism is attributed to local stress depletion in underdense regions, which enhances locally inferred expansion without requiring a global shift of the background expansion parameter.

\[
    \lambda_{\mathrm{MTDF}} = \frac{(1-\beta_{\mathrm{eos}})^2}{1+\alpha}
  \]

\[
    f_{\mathrm{kick}} = \frac{\lambda_{\mathrm{MTDF}}}{24}
  \]

\[
    \frac{\Delta r_s}{r_s} = K_{\mathrm{RS,EFE}} \, f_{\mathrm{kick}} \quad\mathrm{with}\quad K_{\mathrm{RS,EFE}}=-0.215
  \]

Computed from the current V18 fundamental parameters: $\alpha$ = 1.300, $\beta$\textsubscript{eos} = 0.573 gives $\lambda$\textsubscript{MTDF} = 0.079273, f\textsubscript{kick} = 0.003303062, and $\Delta$r\textsubscript{s}/r\textsubscript{s} = -0.000710158. The early field energy fraction f\textsubscript{kick} is also structurally related to the photon-stress coupling parameter $\kappa$ $\approx$ f\textsubscript{kick}/3 (see Companion Part III), linking the gravitational and photon sectors of the theory through a single pair of parameters \{$\alpha$, $\beta$\textsubscript{eos}\}.

\textbf{Status of derived constants.} The factor 1/\allowbreak{}24 in f\textsubscript{kick} = $\lambda$\textsubscript{MTDF}/24 arises from the geometric breakdown $\gamma$\textsubscript{geom} $\cdot$ $\rho$\textsubscript{m}(z\textsubscript{eq})/$\rho$\textsubscript{tot}(z\textsubscript{eq}) $\approx$ (1/\allowbreak{}12) $\times$ (1/\allowbreak{}2); the explicit derivation is given in the short summary paper (\textit{MTDF\_\allowbreak{}short}~[see companion paper], "Hubble solution" section). The sensitivity coefficient K\textsubscript{RS,EFE} = $-$0.215 is an empirical numerical coefficient, calibrated by perturbing the EFE amplitude in the class\_\allowbreak{}mtdf solver and measuring the resulting fractional shift in the sound horizon at decoupling; we report it here as a numerical fit constant, not a first-principles derivation. At the V18 best-fit f\textsubscript{kick} = 3.3 $\times$ 10\textsuperscript{$-$3}, the resulting $\Delta$r\textsubscript{s}/r\textsubscript{s} $\approx$ $-$7 $\times$ 10\textsuperscript{$-$4} is sub-percent and corresponds to a fractional CMB sound-horizon shift well below the Planck statistical uncertainty.

\textbf{Calibration of K\textsubscript{RS,EFE} from a class\_\allowbreak{}mtdf parameter scan.} The coefficient is set by a 9-point scan over the EFE amplitude, varying the dimensionless rescaling parameter k\textsubscript{f} over [0, 0.125, 0.25, 0.5, 0.75, 1.0, 1.5, 2.0, 3.0] (so f\textsubscript{kick}\textsuperscript{eff} = k\textsubscript{f} $\times$ f\textsubscript{kick}\textsuperscript{full} spans 0 to 3 $\times$ the V18 best-fit value). The comoving sound horizon at recombination is read from the class\_\allowbreak{}mtdf thermodynamics module for each run and the fractional shift is taken relative to the k\textsubscript{f} = 0 (no-EFE) baseline:

\begin{table}[H]
\centering
\footnotesize
\begin{tabularx}{\linewidth}{l>{\RaggedRight\arraybackslash}X>{\RaggedRight\arraybackslash}X>{\RaggedRight\arraybackslash}X}
\toprule
\textbf{k\textsubscript{f}} & \textbf{f\textsubscript{kick}\textsuperscript{eff}} & \textbf{r\textsubscript{s} at recombination [Mpc]} & \textbf{$\Delta$r\textsubscript{s}/r\textsubscript{s}} \tabularnewline
\midrule
0.000 & 0.000000 & 150.4109 & 0 \tabularnewline
0.125 & 0.000413 & 150.3971 & $-$9.14 $\times$ 10\textsuperscript{$-$5} \tabularnewline
0.250 & 0.000826 & 150.3834 & $-$1.83 $\times$ 10\textsuperscript{$-$4} \tabularnewline
0.500 & 0.001651 & 150.3558 & $-$3.66 $\times$ 10\textsuperscript{$-$4} \tabularnewline
0.750 & 0.002477 & 150.3283 & $-$5.49 $\times$ 10\textsuperscript{$-$4} \tabularnewline
\textbf{1.000} & \textbf{0.003303} & \textbf{150.3008} & \textbf{$-$7.32 $\times$ 10\textsuperscript{$-$4}} \tabularnewline
1.500 & 0.004954 & 150.2457 & $-$1.10 $\times$ 10\textsuperscript{$-$3} \tabularnewline
2.000 & 0.006606 & 150.1906 & $-$1.46 $\times$ 10\textsuperscript{$-$3} \tabularnewline
3.000 & 0.009909 & 150.0803 & $-$2.20 $\times$ 10\textsuperscript{$-$3} \tabularnewline
\bottomrule
\end{tabularx}
\end{table}

\textbf{Table.} Empirical calibration of the EFE-induced sound-horizon shift. The k\textsubscript{f} = 1 row is the V18 best-fit point; the surrounding rows establish the linear regime over a 30$\times$ range. Cosmological parameters other than k\textsubscript{f} are held at their V18 values (mtdf.ini in the class\_\allowbreak{}mtdf release).A linear least-squares fit $\Delta$r\textsubscript{s}/r\textsubscript{s} = K\textsubscript{RS,EFE} $\cdot$ f\textsubscript{kick} through the origin gives \textbf{K\textsubscript{RS,EFE} = $-$0.222} (RMS residual 4.6 $\times$ 10\textsuperscript{$-$7}); the same fit at the baryon-drag epoch gives K = $-$0.218. Both values agree with the V18 quoted K\textsubscript{RS,EFE} = $-$0.215 within 3 \%, which we take as the residual systematic on this empirical coefficient. A free-intercept fit yields an intercept of 4 $\times$ 10\textsuperscript{$-$7}, statistically indistinguishable from zero, confirming linearity over the scanned range. The relationship is therefore well-approximated as $\Delta$r\textsubscript{s}/r\textsubscript{s} $\approx$ $-$0.22 $\cdot$ f\textsubscript{kick} at recombination, with a numerical sensitivity coefficient that does not require a closed-form derivation but is reproducible from any class\_\allowbreak{}mtdf build by repeating the scan above.

These relations are not presented as a full CMB spectral fit. They exist to provide a transparent mapping from the fundamental parameter set to the specific diagnostic quantity used in the V18 dashboard.

\subsection{Stress tensor energy budget and effective pressure}
At background level, the strict V18 evaluation treats the expansion history as a derived prediction from the fundamental parameter set and the chosen anchors. The underlying physical picture, however, is that accumulated field strain contributes an effective energy density and a negative pressure component that can mimic late time acceleration without invoking a cosmological constant term.

\subsubsection{Core Mechanism}
Cosmic acceleration emerges from negative pressure in the elastic stress tensor $\Sigma$\textsubscript{$\mu$$\nu$} = E S\textsubscript{$\mu$$\nu$}, where spacetime's measurable elastic modulus E = (9.1 $\pm$ 0.4) $\times$ 10$^{-}$$^{1}$$^{0}$ Pa drives expansion through accumulated structural stress.

\textbf{Stress Tensor Energy Density:}

\textbf{u\textsubscript{tensor} = (E/\allowbreak{}2) |S\textsubscript{$\mu$$\nu$}|$^{2}$ = (1/\allowbreak{}2E) $\Sigma$\textsubscript{$\mu$$\nu$} $\Sigma$\textsuperscript{$\mu$$\nu$} [E142]}

\textbf{Negative Pressure:}

\textbf{p\textsubscript{tensor} = -(1/\allowbreak{}3) u\textsubscript{tensor} = -(E/\allowbreak{}6) |S\textsubscript{$\mu$$\nu$}|$^{2}$ [E143]}

\begin{mdframed}[linecolor=accent,linewidth=0.5pt,backgroundcolor=callout-bg,leftline=true,rightline=true,topline=true,bottomline=true,innerleftmargin=8pt,innerrightmargin=8pt,innertopmargin=6pt,innerbottommargin=6pt]
\subsubsection{Dimensional verification:}
[u\textsubscript{tensor}] = [ML$^{-}$$^{1}$T$^{-}$$^{2}$] $\times$ [1]$^{2}$ = [ML$^{-}$$^{1}$T$^{-}$$^{2}$] \checkmark{} (energy density)\newline{}
[p\textsubscript{tensor}] = [ML$^{-}$$^{1}$T$^{-}$$^{2}$] $\times$ [1]$^{2}$ = [ML$^{-}$$^{1}$T$^{-}$$^{2}$] \checkmark{} (pressure)
\end{mdframed}

\textbf{Critical Threshold:} Acceleration begins when stress tensor energy exceeds matter density:

\textbf{u\textsubscript{tensor}(z\textsubscript{t}) > (1/\allowbreak{}3) $\rho$\textsubscript{m} c$^{2}$ at z\textsubscript{t} = 0.74 $\pm$ 0.10 [E144]}

\subsubsection{Modified Friedmann Equation}
\textbf{Stress Tensor Expansion:}

\textbf{H$^{2}$(z) = H$_{0}$$^{2}$[$\Omega$\textsubscript{m}(1+z)$^{3}$ + $\Omega$\textsubscript{tensor} exp(2$\int$$_{0}$\textsuperscript{z} (1+w(z'))/(1+z') dz')] [E145]}

\begin{mdframed}[linecolor=accent,linewidth=0.5pt,backgroundcolor=callout-bg,leftline=true,rightline=true,topline=true,bottomline=true,innerleftmargin=8pt,innerrightmargin=8pt,innertopmargin=6pt,innerbottommargin=6pt]
\subsubsection{Dimensional verification:}
[H$^{2}$] = [T$^{-}$$^{2}$], [$\Omega$ terms] = [1] (dimensionless density parameters) \checkmark{}
\end{mdframed}

\subsubsection{Dynamic Equation of State}
\textbf{Environmental Stress Tensor Coupling:}

\textbf{w(z) = -1 + (2/\allowbreak{}3)$\alpha$(1+z)$^{-}$$^{3}$ - $\beta$\textsubscript{eos} e$^{-}$$^{2}$$^{z}$/$^{z}$$^{t}$ + $\delta$\textsubscript{bf} f\textsubscript{env}(z) [E146]}

\begin{mdframed}[linecolor=accent,linewidth=0.5pt,backgroundcolor=callout-bg,leftline=true,rightline=true,topline=true,bottomline=true,innerleftmargin=8pt,innerrightmargin=8pt,innertopmargin=6pt,innerbottommargin=6pt]
\subsubsection{Status of the environmental coupling term}
The raw MTDF equation of state, evaluated without environmental coupling, predicts w$_{0}$ = -0.703, in 4.1$\sigma$ tension with current constraints. We report this failure rather than omitting it. The environmental term $\delta$\textsubscript{bf} f\textsubscript{env}(z) in E146 was introduced after this tension was identified, and we state plainly that it is a post-hoc extension of the equation of state. Its amplitude, however, is not a free rescue parameter: $\delta$\textsubscript{bf} = 0.150 $\pm$ 0.030 is fixed externally by cluster-scale baryon mass fractions measured independently of any cosmological fit (Giodini et al. 2009; Gonzalez et al. 2013). With that externally fixed value the prediction becomes w$_{0}$ = -1.013 $\pm$ 0.025, in 0.4$\sigma$ agreement with Planck + SH0ES.

We regard the absence of a first-principles derivation of $\delta$\textsubscript{bf}, and of the unit coefficient with which it enters E146, as the largest open theoretical weakness of the framework (Section 8). It also supplies a falsification channel: because $\delta$\textsubscript{bf} is chained to cluster baryon accounting, future measurements placing the intermediate-mass cluster baryon fraction outside approximately 0.12 to 0.18 would degrade the equation-of-state prediction, and the framework forbids re-tuning $\delta$\textsubscript{bf} to compensate.
\end{mdframed}

\subsection{Void stress depletion and late time expansion}
In the dual epoch narrative used throughout the framework, early time effects are separated from late time effects. For the late time component, MTDF motivates a void-driven depletion mechanism in which stress is redistributed as structure grows, leading to a scale dependent modification of the effective expansion rate.

\subsubsection{Local Hubble Enhancement}
Reduced stress tensor magnitude in cosmic voids modifies local expansion through environmental coupling:

\textbf{H\textsubscript{local}/H$_{0}$ = 1 + ($\alpha$\textsubscript{void}/2) tanh(r/\allowbreak{}$\beta$) = 1.084 $\pm$ 0.015 [E147]} {\small (evaluated at KBC-level void depth; conditional, see callout below)}

\begin{mdframed}[linecolor=accent,linewidth=0.5pt,backgroundcolor=callout-bg,leftline=true,rightline=true,topline=true,bottomline=true,innerleftmargin=8pt,innerrightmargin=8pt,innertopmargin=6pt,innerbottommargin=6pt]
\subsubsection{Conditional local-environment consistency}
MTDF is consistent with a local-environment route to the Hubble-tension problem: a sufficiently deep local underdensity raises the locally inferred expansion rate relative to the global background value. For a KBC/\allowbreak{}Haslbauer-level contrast ($\delta$ $\approx$ 0.46), the resulting enhancement can reach H\textsubscript{local} $\approx$ 73.1 km/\allowbreak{}s/\allowbreak{}Mpc, consistent with SH0ES (73.04 $\pm$ 1.04 km/\allowbreak{}s/\allowbreak{}Mpc; Riess et al. 2022). For the shallower contrast favoured by cluster-based reconstructions (B\"ohringer et al. 2020, $\delta$ $\approx$ 0.2 to 0.3), the same mechanism gives only partial alleviation, H\textsubscript{local} $\approx$ 70 to 71 km/\allowbreak{}s/\allowbreak{}Mpc.

The present strict formulation therefore classifies H\textsubscript{local} = 73.1 km/\allowbreak{}s/\allowbreak{}Mpc as a conditional local-environment consistency rather than an independently derived prediction: its viability depends on the still-contested depth and geometry of the local underdensity over the relevant distance-ladder volume. The relation S\textsubscript{0} = $\alpha$$\sqrt{\,}$$\Omega$\textsubscript{DE} is an algebraic identity following from the definition of the elastic modulus E, and is not treated as the physical mechanism producing the local expansion value. The decisive remaining tests are whether MTDF's growth history makes a KBC-depth local environment more or less probable than in $\Lambda$CDM, and whether the predicted low-redshift environmental transition exhibits the expected directional (anisotropic) dependence.

\textbf{Note:} The base value 67.4 km/\allowbreak{}s/\allowbreak{}Mpc is the Planck $\Lambda$CDM background posterior. The workbook anchor H$_{0}$ = 70.0 km/\allowbreak{}s/\allowbreak{}Mpc (Table 2) is a calibration choice for the validation pipeline distance calculations and is not used here.
\end{mdframed}

\textbf{Mechanism:} Reduced stress tensor in KBC void:

\textbf{$\|$$\Sigma$\textsuperscript{void}\textsubscript{$\mu$$\nu$}$\|$ = 0.75 $\|$$\Sigma$\textsuperscript{cosmic}\textsubscript{$\mu$$\nu$}$\|$ [E148]}

\begin{mdframed}[linecolor=accent,linewidth=0.5pt,backgroundcolor=callout-bg,leftline=true,rightline=true,topline=true,bottomline=true,innerleftmargin=8pt,innerrightmargin=8pt,innertopmargin=6pt,innerbottommargin=6pt]
\subsubsection{Dimensional verification:}
[$\|$$\Sigma$\textsuperscript{void}\textsubscript{$\mu$$\nu$}$\|$/$\|$$\Sigma$\textsuperscript{cosmic}\textsubscript{$\mu$$\nu$}$\|$] = [ML$^{-}$$^{1}$T$^{-}$$^{2}$]/[ML$^{-}$$^{1}$T$^{-}$$^{2}$] = [1] \checkmark{} (stress tensor ratio)
\end{mdframed}

\subsection{Stress evolution and near term forecasts}
Because the stress field is dynamical and relaxes over a finite timescale, the model implies not only a present day expansion rate but also an evolution that can be tested by future probes. In V18, these appear as falsifiable targets rather than as calibration inputs.

\subsubsection{Tensor Energy Density Evolution}
\textbf{Stress Tensor Energy:}

\textbf{u\textsubscript{tensor}(z) = u$_{0}$[1 + $\alpha$ ln(1+z)] e$^{-}$$\gamma$t(z) [E149]}

\begin{mdframed}[linecolor=accent,linewidth=0.5pt,backgroundcolor=callout-bg,leftline=true,rightline=true,topline=true,bottomline=true,innerleftmargin=8pt,innerrightmargin=8pt,innertopmargin=6pt,innerbottommargin=6pt]
\subsubsection{Dimensional verification:}
[u$_{0}$] = [ML$^{-}$$^{1}$T$^{-}$$^{2}$] (reference energy density)\newline{}
[$\alpha$ ln(1+z)] = [1] $\times$ [1] = [1] \checkmark{}\newline{}
[e$^{-}$$\gamma$t(z)] = [1] \checkmark{}
\end{mdframed}

where u$_{0}$ = (E/\allowbreak{}2)$\alpha$$^{2}$ = 7.7 $\times$ 10$^{-}$$^{1}$$^{0}$ J/\allowbreak{}m$^{3}$ sets the characteristic stress tensor energy scale.

\subsubsection{Future Evolution Timeline}
\textbf{Testable Predictions:}

\begin{itemize}
  \item \textbf{Acceleration parameter:} q$_{0}$ = -0.55 $\rightarrow$ -0.45 $\rightarrow$ -0.30
  \item \textbf{Observable signature:} Gradual deceleration of cosmic acceleration
  \item \textbf{Critical transition:} Clear w(z) evolution testable with future surveys
\end{itemize}

\section{Emergent phenomenology across scales}
\label{sec:sec4}
MTDF is intended to unify several empirical patterns: flat rotation curves, the radial acceleration relation, galaxy regularities, cluster scale offsets, and the large scale expansion history. This master paper separates (i) the conceptual mechanism that MTDF proposes and (ii) the strict numerical validation scope currently implemented.

\subsection{Galaxy scale dynamics}
On galaxy scales, the framework attributes the regularity of rotation curves to the response of the field to baryonic distributions and to the strain-lever scale. This is reflected in scalar benchmarks such as the rotation curve scatter and the intrinsic scatter of the radial acceleration relation. In the strict protocol these are reported as benchmarks rather than as independent validation targets, because parts of the mapping are used to define internal consistency constraints for the framework.

\subsubsection{Mechanism: stress divergence as an additional acceleration channel}
On galaxy scales, MTDF explains the near universality of rotation curves as an emergent response of the field to baryonic structure in the low acceleration regime. In the non relativistic limit, the additional contribution may be expressed schematically as a stress divergence term added to the Newtonian acceleration:

\[
\frac{d\mathbf{v}_*}{dt} = -\nabla\Phi_\mathrm{Newt} - \frac{1}{\rho_*}\,\nabla_j F^j{}_t
\]

\subsubsection{Physical Foundation}
Outer stars experience deceleration due to spacetime stress divergence, governed by the elastic stress tensor $\Sigma$\textsubscript{$\mu$$\nu$} = E S\textsubscript{$\mu$$\nu$}:

\textbf{dv\textsubscript{*}/dt = -$\nabla$$\Phi$\textsubscript{Newt} - (1/\allowbreak{}$\rho$\textsubscript{*})$\nabla$\textsubscript{j}$\Sigma$\textsuperscript{j}\textsubscript{t} [E16]}

where $\Sigma$\textsuperscript{j}\textsubscript{t} = time-space components of stress tensor representing the resistance force density from spacetime deformation.

\begin{mdframed}[linecolor=accent,linewidth=0.5pt,backgroundcolor=callout-bg,leftline=true,rightline=true,topline=true,bottomline=true,innerleftmargin=8pt,innerrightmargin=8pt,innertopmargin=6pt,innerbottommargin=6pt]
\subsubsection{Dimensional verification:}
[$\nabla$$\Phi$] = [L$^{2}$T$^{-}$$^{2}$]/[L] = [LT$^{-}$$^{2}$] (gravitational acceleration)\newline{}
[(1/\allowbreak{}$\rho$)$\nabla$\textsubscript{j}$\Sigma$\textsuperscript{j}\textsubscript{t}] = [M$^{-}$$^{1}$L$^{3}$] $\times$ [L$^{-}$$^{1}$][ML$^{-}$$^{1}$T$^{-}$$^{2}$] = [LT$^{-}$$^{2}$] (stress acceleration)\newline{}
Total: [LT$^{-}$$^{2}$] = [LT$^{-}$$^{2}$] \checkmark{}
\end{mdframed}

\subsubsection{Critical Acceleration Threshold}
The stress tensor effects become dominant when stress divergence approaches gravitational acceleration:

\textbf{a\textsubscript{crit} = $\alpha$c$^{2}$/$\beta$ = (1.67 $\pm$ 0.33) $\times$ 10$^{-}$$^{7}$ m/\allowbreak{}s$^{2}$ [E17]}

\begin{mdframed}[linecolor=accent,linewidth=0.5pt,backgroundcolor=callout-bg,leftline=true,rightline=true,topline=true,bottomline=true,innerleftmargin=8pt,innerrightmargin=8pt,innertopmargin=6pt,innerbottommargin=6pt]
\subsubsection{Dimensional verification:}
[$\alpha$c$^{2}$/$\beta$] = [1] $\times$ [L$^{2}$T$^{-}$$^{2}$] / [L] = [LT$^{-}$$^{2}$] \checkmark{} (acceleration units)
\end{mdframed}

This scale characterises where stress tensor contributions to the effective potential become comparable to Newtonian gravity within the historical effective parameterisation; it is not a sharp threshold, and the effective phenomenology uses the full enhancement factor 1 + $\alpha$/(1 + r/\allowbreak{}$\beta$). Note that a\textsubscript{crit} $\approx$ 10$^{-}$$^{7}$ m/\allowbreak{}s$^{2}$ is a different scale from the empirical acceleration scale a\textsubscript{0} $\approx$ 1.2 $\times$ 10$^{-}$$^{1}$$^{0}$ m/\allowbreak{}s$^{2}$. In the frozen unified action (Section 2A) the genuine acceleration scale of the galactic response is the solid-anchored a\textsubscript{MTDF} = $\kappa$\textsubscript{a}$\sqrt{\,}$(G u\textsubscript{bg}), which sits at the a\textsubscript{0} scale up to the O(1) constant $\kappa$\textsubscript{a}: the canonical low-Q branch is a true low-acceleration (aquadratic) response, and the constant-enhancement description below is its historical effective limit.

\begin{mdframed}[linecolor=accent,linewidth=0.5pt,backgroundcolor=callout-bg,leftline=true,rightline=true,topline=true,bottomline=true,innerleftmargin=8pt,innerrightmargin=8pt,innertopmargin=6pt,innerbottommargin=6pt]
\subsubsection{Critical Acceleration Parameter}
\textbf{Tensor Origin:} Derived from stress tensor divergence threshold where $\nabla$\textsubscript{j}$\Sigma$\textsuperscript{j}\textsubscript{t}/$\rho$ $\approx$ $\nabla$$\Phi$
\end{mdframed}

\subsubsection{Tensor-based rotation curve solution and the strain-lever scale}

\subsubsection{Modified Potential}
In the weak-field limit, MTDF modifies the effective gravitational potential through spacetime strain effects derived from the stress tensor:

\textbf{$\Phi$\textsubscript{eff} = $\Phi$\textsubscript{Newt} - (E/\allowbreak{}$\rho$\textsubscript{tensor}) $\|$S\textsubscript{$\mu$$\nu$}$\|$ ln(1+r/\allowbreak{}$\beta$) [E18]}

where the strain contribution emerges from the dimensionless strain tensor S\textsubscript{$\mu$$\nu$} = $\beta$ $\nabla$\textsubscript{$\mu$}$\xi$\textsubscript{$\nu$}, with E = 9.1$\times$10$^{-}$$^{1}$$^{0}$ Pa representing the elastic modulus of spacetime.

\subsubsection{Velocity Profile}
The circular velocity becomes:

\textbf{v\textsubscript{c}$^{2}$(r) = GM/\allowbreak{}r $\times$ [1 + $\alpha$/(1 + r/\allowbreak{}$\beta$)] [E20]}

\begin{mdframed}[linecolor=accent,linewidth=0.5pt,backgroundcolor=callout-bg,leftline=true,rightline=true,topline=true,bottomline=true,innerleftmargin=8pt,innerrightmargin=8pt,innertopmargin=6pt,innerbottommargin=6pt]
\subsubsection{Dimensional verification:}
\begin{itemize}
  \item GM/\allowbreak{}r: [L$^{3}$T$^{-}$$^{2}$]/[L] = [L$^{2}$T$^{-}$$^{2}$] \checkmark{}
  \item $\alpha$/(1 + r/\allowbreak{}$\beta$): [1]/[1] = [1] \checkmark{}
  \item Total: [L$^{2}$T$^{-}$$^{2}$] $\rightarrow$ v\textsubscript{c} in [LT$^{-}$$^{1}$] \checkmark{}
\end{itemize}
\end{mdframed}

\textbf{Tensor origin:} $\alpha$ $\propto$ |S\textsubscript{$\mu$$\nu$}|\textsubscript{cosmic} represents the characteristic strain magnitude from cosmic stress tensor dynamics.

\subsubsection{Regime behaviour}
At galactic scales (r $\ll$ $\beta$ = 22,685 kpc), the enhancement factor is approximately constant: v\textsubscript{c}$^{2}$(r) $\approx$ (1+$\alpha$) GM(r)/r, a uniform 2.3$\times$ boost to the Newtonian prediction. Flat rotation curves arise from this constant enhancement combined with the extended baryonic mass distribution M(r). At scales r $\gg$ $\beta$, the enhancement vanishes and standard Keplerian behaviour is recovered.

\subsubsection{Radial acceleration relation and scatter control}
The MTDF enhancement enters the radial acceleration relation through the stress tensor coupling:

\textbf{g\textsubscript{obs} = g\textsubscript{bar} $\times$ [1 + $\alpha$/(1 + r/\allowbreak{}$\beta$)] [E22]}

\begin{mdframed}[linecolor=accent,linewidth=0.5pt,backgroundcolor=callout-bg,leftline=true,rightline=true,topline=true,bottomline=true,innerleftmargin=8pt,innerrightmargin=8pt,innertopmargin=6pt,innerbottommargin=6pt]
\subsubsection{Radial acceleration relation: canonical resolution of the former open issue}
At the radii probed by SPARC (r $\ll$ $\beta$) the factor in E22 is approximately constant, so as written it shifts the relation by a roughly uniform offset rather than producing the curved, slope-1/\allowbreak{}2 bend observed at low accelerations. Earlier versions of this paper recorded that shape-level limitation as an open problem. The frozen unified action resolves it: the low-Q branch of the constitutive J-function is the aquadratic (A0 cubic) response, giving the deep-MOND relation g = $\sqrt{\,}$(g\textsubscript{bar} a\textsubscript{MTDF}) and hence the observed low-acceleration curvature and the M\textsuperscript{1/\allowbreak{}4} baryonic Tully-Fisher scaling; the radial acceleration relation is validated against SPARC at \~{}0.11 dex scatter. The epistemic status is stated exactly: the aquadratic low-Q form is a declared postulate (honest form, with the RAR/\allowbreak{}BTFR as forcing ancestry), the anchor relation a\textsubscript{MTDF} = $\kappa$\textsubscript{a}$\sqrt{\,}$(G u\textsubscript{bg}) is derived, and the RAR confrontation itself is validated. The relation retains no assembly dependence at galaxy scale (universal RAR, validated); environmental dependence at fixed baryonic structure survives only as the conditional prediction discussed in Sections 2A.2 and 10.4.
\end{mdframed}

\begin{mdframed}[linecolor=accent,linewidth=0.5pt,backgroundcolor=callout-bg,leftline=true,rightline=true,topline=true,bottomline=true,innerleftmargin=8pt,innerrightmargin=8pt,innertopmargin=6pt,innerbottommargin=6pt]
\subsubsection{Section 4 summary}
\begin{itemize}
  \item \textbf{P1-SPARC:} benchmark pillar; achieved RMS residual scatter 0.185 dex against a target of 0.174 $\pm$ 0.011 dex (both below the 0.25 dex reference threshold). This is an aggregate scatter statistic, not a per-galaxy shape fit (see below).
  \item \textbf{Parameters:} $\alpha$ = 1.30, $\beta$ = 22.7 Mpc, frozen and void-motivated rather than independently measured (Section 5.0).
  \item \textbf{Scope:} the same tensor formalism is applied at galactic and cluster scales; galaxy-scale pillars are treated as scalar benchmarks and are excluded from the strict scoring totals.
\end{itemize}
\end{mdframed}

This validation suite supports MTDF as a quantitatively predictive framework for galactic dynamics through stress tensor mechanics, under the stated test protocol, providing a basis for its extension to larger cosmic scales within the present validation framework.

\begin{mdframed}[linecolor=accent,linewidth=0.5pt,backgroundcolor=callout-bg,leftline=true,rightline=true,topline=true,bottomline=true,innerleftmargin=8pt,innerrightmargin=8pt,innertopmargin=6pt,innerbottommargin=6pt]
\subsubsection{How SPARC appears in the V18 strict protocol}
In the V18 strict protocol, SPARC is a scalar benchmark, not a direct stress tensor validation: parts of the galaxy scale mapping are treated as internal consistency constraints and scalar benchmarks, while the strict totals emphasise independent vector likelihood targets. This preserves falsifiability while avoiding circular use of the same information as both calibration and validation.

\subsubsection{Galaxy-sector statistic and its scope}
\begin{mdframed}[linecolor=accent,linewidth=0.5pt,backgroundcolor=callout-bg,leftline=true,rightline=true,topline=true,bottomline=true,innerleftmargin=8pt,innerrightmargin=8pt,innertopmargin=6pt,innerbottommargin=6pt]
\textbf{What the statistic is.} On the SPARC sample (Lelli et al. 2016, AJ 152, 157), the enhancement law reproduces the rotation-curve sample with an achieved RMS residual scatter of RMS\textsubscript{log} = 0.185 dex in log-log space across 175 galaxies (target 0.174 $\pm$ 0.011 dex; a 5.4\% reduction from the Newtonian value of 0.195 dex), below the 0.25 dex reference threshold. This is an aggregate scatter of the residual log(v\textsubscript{obs}/v\textsubscript{model}); it is not a per-galaxy rotation-curve shape fit, and it is evaluated with the mass-to-light treatment carried in the SPARC mass models rather than a free per-galaxy normalisation.

\textbf{Shape versus amplitude.} At radii far inside the strain-lever scale $\beta$ the enhancement factor is approximately constant at (1 + $\alpha$) = 2.3. A constant enhancement rescales the amplitude of a rotation curve without altering its shape, so within this effective parameterisation the reproduction of rotation-curve flatness depends on the baryonic mass profile M(r) and the mass-to-light treatment. In the canonical architecture the shape-level question is closed by the low-Q aquadratic branch (Section 2A.1); the statistic reported here remains an amplitude-level benchmark of the historical effective law.

\textbf{Scope.} Because of this, the galaxy-scale pillars are treated as scalar benchmarks and are excluded from the strict scoring totals, so this result does not contaminate the headline metric. The median stress-induced velocity contribution is 12.7 km/\allowbreak{}s.
\end{mdframed}

\subsubsection{Outliers and anomaly handling}
MTDF includes a qualitative framework for interpreting persistent outliers (for example, strong non equilibrium systems, interacting galaxies, or environments with unusual baryonic distributions). In V18 these are treated as diagnostic cases and not used to tune parameters.

\subsubsection{Interpretive framework for outliers}
\begin{mdframed}[linecolor=accent,linewidth=0.5pt,backgroundcolor=callout-bg,leftline=true,rightline=true,topline=true,bottomline=true,innerleftmargin=8pt,innerrightmargin=8pt,innertopmargin=6pt,innerbottommargin=6pt]
\subsubsection{Tensor-Based Anomaly Resolution}
\textbf{Mechanistic Foundation:} Outlier interpretations are organised around the stress tensor divergence $\nabla$\textsubscript{$\mu$}$\Sigma$\textsuperscript{$\mu$$\nu$} = 8$\pi$G/\allowbreak{}c$^{4}$ T\textsubscript{$\mu$$\nu$}, combined with environment-dependent coupling factors where needed
\end{mdframed}

\begin{mdframed}[linecolor=accent,linewidth=0.5pt,backgroundcolor=callout-bg,leftline=true,rightline=true,topline=true,bottomline=true,innerleftmargin=8pt,innerrightmargin=8pt,innertopmargin=6pt,innerbottommargin=6pt]
\subsubsection{Environmental Parameter Derivation}
\textbf{Environmental Coupling Relationship:} $\alpha$\textsubscript{cluster} = 0.7$\alpha$\textsubscript{void}

\textbf{Derivation Reference:} This relationship emerges from the strain tensor saturation model in high-density environments, where $\gamma$\textsubscript{cut} = $\beta$\textsubscript{eos}$^{3}$/$\alpha$ = 0.144 limits stress tensor-matter coupling. The factor 0.7 represents the average suppression factor observed in cluster environments compared to void regions, as documented in the V18 validation workbook.

\textbf{Physical Interpretation:} In high-density cluster environments, the increased matter density reduces the effective stress tensor-matter coupling strength by approximately 30\% compared to low-density void regions, naturally explaining the reduced stress tensor effects in cluster satellite systems.
\end{mdframed}
\end{mdframed}

\subsection{Structure growth and large scale surveys}
Large scale structure is tested via vector likelihoods that encode both distances and growth. The current strict totals include Pantheon+ SNe, DESI Y1 BAO, cosmic chronometer H(z), and DR16 f$\sigma$$_{8}$ growth. These tests jointly probe the background expansion and the normalisation of growth, with the growth amplitude fixed at the sealed value (the former fitted-nuisance treatment is retired; the DR16 entry is carried as PREDICTION-CONTESTED).

\subsubsection{Stress tensor properties relevant for structure growth}
Beyond galaxy scales, MTDF treats the strain field as a coherent medium whose gradients and relaxation dynamics shape the emergence of filaments, voids, and scale dependent clustering. The following summaries provide the qualitative mechanism that motivates the large scale predictions evaluated in the strict V18 artefacts.

The spacetime stress tensor possesses the following characteristics based on its mathematical description through the dimensionless strain tensor S\textsubscript{$\mu$$\nu$} = $\beta$ $\nabla$\textsubscript{$\mu$} $\xi$\textsubscript{$\nu$}:

\begin{itemize}
  \item \textbf{Pervasive Nature:} Unlike the traditional view of space as empty vacuum, this stress tensor provides a structured medium through which all mass and energy move, characterized by elastic modulus E = (9.1 $\pm$ 0.4) $\times$ 10$^{-}$$^{1}$$^{0}$ Pa.
  \item \textbf{Weak Interaction:} Stress tensor effects are subtle, unobservable on small timescales or distances, but cumulatively significant over cosmic scales due to the ultra-small relaxation rate $\gamma$ = 2.437 $\times$ 10$^{-}$$^{1}$$^{8}$ s$^{-}$$^{1}$.
  \item \textbf{Stress Storage:} Objects moving through spacetime experience minute resistance, with the lost kinetic energy stored as stress tensor magnitude with energy density u\textsubscript{tensor} = (E/\allowbreak{}2)|S\textsubscript{$\mu$$\nu$}|$^{2}$ rather than being destroyed.
  \item \textbf{Dynamic Response:} The stress tensor actively responds to mass concentrations, creating stress gradients and tensor structures around massive objects according to the equation G\textsuperscript{$\mu$$\nu$}[S] = (8$\pi$G/\allowbreak{}c$^{4}$) T\textsuperscript{$\mu$$\nu$}.
\end{itemize}

The spacetime stress tensor remains active at all scales, but its effects become dominant only in low-acceleration environments. In regions of strong gravitational curvature or rapid motion (e.g. near stellar masses), stress tensor gradients are suppressed. This behaviour can be characterised by a natural relaxation timescale:

\textbf{$\tau$\textsubscript{relax} = $\chi$ / a$^{2}$ [E2]}

\begin{mdframed}[linecolor=accent,linewidth=0.5pt,backgroundcolor=callout-bg,leftline=true,rightline=true,topline=true,bottomline=true,innerleftmargin=8pt,innerrightmargin=8pt,innertopmargin=6pt,innerbottommargin=6pt]
\subsubsection{Stress Tensor Relaxation Parameter}
\textbf{Derivation:} Stress tensor compliance $\chi$ = $\eta$ / $\rho$\textsubscript{tensor}, where $\eta$ \~{} E$\cdot$$\tau$ $\approx$ 3.7 $\times$ 10$^{8}$ Pa$\cdot$s is an order-of-magnitude stress tensor viscosity and u$_{0}$ = 7.7 $\times$ 10$^{-}$$^{1}$$^{0}$ J/\allowbreak{}m$^{3}$ is the characteristic stress tensor energy density (= (E/\allowbreak{}2)$\alpha$$^{2}$)
\end{mdframed}

\begin{mdframed}[linecolor=accent,linewidth=0.5pt,backgroundcolor=callout-bg,leftline=true,rightline=true,topline=true,bottomline=true,innerleftmargin=8pt,innerrightmargin=8pt,innertopmargin=6pt,innerbottommargin=6pt]
\subsubsection{Dimensional verification:}
[$\chi$] = [T$^{3}$L$^{-}$$^{1}$] (stress tensor compliance), [a$^{2}$] = [L$^{2}$T$^{-}$$^{4}$]\newline{}
[$\chi$/a$^{2}$] = [T$^{3}$L$^{-}$$^{1}$]/[L$^{2}$T$^{-}$$^{4}$] = [T] \checkmark{} (time units)
\end{mdframed}

where \textbf{$\chi$} represents the stress tensor's compliance with dimensions [T$^{3}$L$^{-}$$^{1}$], and \textbf{a} is the local acceleration. High-acceleration environments inhibit the coherent buildup of stress tensor magnitude, rendering the tensor effectively negligible in regions where General Relativity already provides accurate predictions.

\subsubsection{Stress accumulation, relaxation cycles, and hierarchy}
Building on the elastic stress tensor framework (Section 2.3), spacetime operates through fundamental stress-relaxation cycles governed by the viscoelastic evolution equation:

\textbf{$\partial$$^{2}$$\phi$/$\partial$t$^{2}$ - c$^{2}$$\nabla$$^{2}$$\phi$ + $\gamma$$\partial$$\phi$/$\partial$t = 8$\pi$G$\rho$ [E1]}

\subsubsection{Stress Buildup Phase}
The elastic stress tensor $\Sigma$\textsubscript{$\mu$$\nu$} = E S\textsubscript{$\mu$$\nu$} accumulates energy through multiple mechanisms:

\begin{itemize}
  \item \textbf{Galactic Motion:} As galaxies rotate and evolve, their movement creates increasing stress in surrounding spacetime, characterized by strain magnitude |S\textsubscript{$\mu$$\nu$}| \~{} $\alpha$ = 1.30
  \item \textbf{Stellar Orbits:} Individual stars experience ultra-slow deceleration (approximately 0.001\% per million years) as they interact with stress tensor resistance over the characteristic time $\tau$ = 13.0 $\pm$ 0.2 Gyr
  \item \textbf{Energy Storage:} Lost kinetic energy is stored as increasing stress tensor magnitude with energy density u\textsubscript{tensor} = (E/\allowbreak{}2)|S\textsubscript{$\mu$$\nu$}|$^{2}$ rather than being dissipated
  \item \textbf{Critical Accumulation:} Stress builds over billions of years until reaching critical thresholds defined by the strain-lever length scale $\beta$ = 22.7 Mpc
\end{itemize}

\subsubsection{Relaxation Phase}
When stress tensor gradients exceed critical thresholds, the elastic spacetime undergoes structured energy release:

\begin{itemize}
  \item \textbf{High-Energy Events:} When stress exceeds critical levels |$\nabla$$\Sigma$\textsubscript{$\mu$$\nu$}| > T\textsubscript{crit} = 2.1 $\times$ 10$^{-}$$^{9}$ Pa, sudden releases occur through quasars, black holes, and other high-energy phenomena
  \item \textbf{Cosmic Expansion:} Large-scale stress tensor relaxation drives cosmic expansion through negative pressure p\textsubscript{tensor} = -(1/\allowbreak{}3)u\textsubscript{tensor} as the universe attempts to reduce overall stress
  \item \textbf{Structural Reorganization:} Elastic stress tensor gradients reconfigure, creating new cosmic web patterns over the relaxation time $\tau$
  \item \textbf{Energy Redistribution:} Released energy redistributes throughout the stress tensor system while maintaining total energy conservation
\end{itemize}

\subsubsection{Cyclical Nature}
This process operates continuously across multiple timescales-from galactic rotation cycles (hundreds of millions of years) to cosmic evolution cycles (billions of years), creating the dynamic universe we observe. The process is characterized by the universal equation of state:

\textbf{w(z) = -1 + (2/\allowbreak{}3)$\alpha$(1+z)$^{-}$$^{3}$ - $\beta$\textsubscript{eos} e$^{-}$$^{2}$$^{z}$/$^{z}$$^{t}$ + $\delta$\textsubscript{bf} f\textsubscript{env}(z) [E3]}

where $\beta$\textsubscript{eos} = 0.573 $\pm$ 0.012 emerges from universal critical amplitudes. This is a phenomenological MTDF fit; the physics-analogue DOI (HotQCD 2012) is not the source of the value.

\subsubsection{Topology, defects, and filament formation}
The elastic stress tensor $\Sigma$\textsubscript{$\mu$$\nu$} = E S\textsubscript{$\mu$$\nu$} creates hierarchical structure through scale-dependent interactions:

\subsubsection{Multi-Scale Structure Formation}
\begin{itemize}
  \item \textbf{Around Individual Masses:} Elastic stress tensor gradients wrap around planets, stars, and galaxies, creating localized stress patterns that govern orbital mechanics and stellar dynamics with characteristic strain |S\textsubscript{$\mu$$\nu$}| \~{} $\alpha$
  \item \textbf{Galactic Scale:} Stress tensor gradients encompass entire galaxies, with stress tensor magnitude building as stars orbit through spacetime over billions of years, creating the observed flat rotation curves through the modified potential $\Phi$\textsubscript{eff}(r) = $\Phi$\textsubscript{Newton}(r) + $\Phi$\textsubscript{strain}(r)
  \item \textbf{Cosmic Scale:} Stress tensor gradients extend throughout the universe, creating a web-like structure. Galaxy clusters act as gravitational "poles," generating stress tensor gradients that span millions of light-years with the strain-lever length scale $\beta$ = 22.7 Mpc
\end{itemize}

The stress tensor's properties naturally create the observed cosmic web through elastic dynamics, where matter concentrates along stress tensor gradients creating dense filaments, while regions between gradients form vast cosmic voids.

\begin{mdframed}[linecolor=accent,linewidth=0.5pt,backgroundcolor=callout-bg,leftline=true,rightline=true,topline=true,bottomline=true,innerleftmargin=8pt,innerrightmargin=8pt,innertopmargin=6pt,innerbottommargin=6pt]
\subsubsection{Stress Tensor Structure Formation}
The elastic stress tensor $\Sigma$\textsubscript{$\mu$$\nu$} organizes matter through:

\begin{itemize}
  \item \textbf{Stress Concentration:} |$\Sigma$\textsubscript{$\mu$$\nu$}| \~{} E$\alpha$ = (9.1$\times$10$^{-}$$^{1}$$^{0}$)(1.30) = 1.2$\times$10$^{-}$$^{9}$ Pa at galactic centers
  \item \textbf{Stress Gradients:} $\nabla$$\Sigma$\textsubscript{$\mu$$\nu$} drives matter flow along minimum resistance paths
  \item \textbf{Strain-lever scale:} $\beta$ = 22.7 Mpc is the fundamental strain-lever length (\(\tilde{S} = \beta\,\nabla\xi\)); it sets the range of the active stress response, not a correlation length of a fluctuating field
  \item \textbf{Energy Density:} u = (E/\allowbreak{}2)|S\textsubscript{$\mu$$\nu$}|$^{2}$ = 7.7$\times$10$^{-}$$^{1}$$^{0}$ J/\allowbreak{}m$^{3}$ characteristic stress tensor energy
\end{itemize}

\textbf{Mathematical Foundation:} $\Sigma$\textsubscript{$\mu$$\nu$} = E S\textsubscript{$\mu$$\nu$} = E $\beta$ $\nabla$\textsubscript{$\mu$}$\xi$\textsubscript{$\nu$} converts dimensionless strain to measurable stress
\end{mdframed}

\subsubsection{Cosmic web topology and causal propagation}

\subsubsection{Cosmic Structure Mapping}
The stress tensor $\Sigma$\textsubscript{$\mu$$\nu$} organizes matter via:

\textbf{d$^{2}$x$^{i}$/dt$^{2}$ = -$\Gamma$$^{i}$\textsubscript{jk}$\dot{x}$$^{j}$$\dot{x}$$^{k}$ + (1/\allowbreak{}$\rho$)$\nabla$\textsubscript{j}$\Sigma$\textsuperscript{i$_{j}$} [E6']}

\subsubsection{Structure Formation Mechanics}
Matter flow velocity:

\textbf{v\textsubscript{flow}$^{i}$ = -(D\textsubscript{tensor}/$\rho$) g$^{i}$$^{j}$ $\nabla$\textsubscript{j}$\|$$\Sigma$\textsubscript{$\mu$$\nu$}$\|$ [E8]}

where D\textsubscript{tensor} = E/\allowbreak{}($\rho$\textsubscript{tensor} $\gamma$) = 4.85 $\times$ 10$^{1}$$^{7}$ m$^{2}$/s

\begin{mdframed}[linecolor=accent,linewidth=0.5pt,backgroundcolor=callout-bg,leftline=true,rightline=true,topline=true,bottomline=true,innerleftmargin=8pt,innerrightmargin=8pt,innertopmargin=6pt,innerbottommargin=6pt]
\subsubsection{Dimensional verification:}
[D/\allowbreak{}$\rho$] = [L$^{2}$T$^{-}$$^{1}$]/[ML$^{-}$$^{3}$] = [ML$^{-}$$^{1}$L$^{2}$T$^{-}$$^{1}$]

[$\nabla$$\|$F$\|$] = [L$^{-}$$^{1}$][ML$^{-}$$^{1}$T$^{-}$$^{2}$] = [ML$^{-}$$^{2}$T$^{-}$$^{2}$L$^{-}$$^{1}$]

[v] = [ML$^{-}$$^{1}$L$^{2}$T$^{-}$$^{1}$] $\times$ [ML$^{-}$$^{2}$T$^{-}$$^{2}$L$^{-}$$^{1}$] = [LT$^{-}$$^{1}$] \checkmark{}
\end{mdframed}

\subsubsection{Tensor Equation with Causality}
\textbf{$\square$$\Sigma$\textsubscript{$\mu$$\nu$} - $\gamma$$\partial$\textsubscript{t}$\Sigma$\textsubscript{$\mu$$\nu$} = (8$\pi$G/\allowbreak{}c$^{4}$)T\textsubscript{$\mu$$\nu$}\textsuperscript{matter} [E1']}

where $\square$ = $\partial$$_{t}$$^{2}$ - c$^{2}$$\nabla$$^{2}$ ensures c-limited propagation

\begin{mdframed}[linecolor=accent,linewidth=0.5pt,backgroundcolor=callout-bg,leftline=true,rightline=true,topline=true,bottomline=true,innerleftmargin=8pt,innerrightmargin=8pt,innertopmargin=6pt,innerbottommargin=6pt]
\subsubsection{Material Parameters Table}
\textbf{Reference:} Complete parameter derivations are summarised in Section 5 and cross-referenced to the validation workbook.
\end{mdframed}

\begin{mdframed}[linecolor=accent,linewidth=0.5pt,backgroundcolor=callout-bg,leftline=true,rightline=true,topline=true,bottomline=true,innerleftmargin=8pt,innerrightmargin=8pt,innertopmargin=6pt,innerbottommargin=6pt]
\subsubsection{Scope note}
The material in this subsection is theory-motivated mechanism building. Quantitative claims that are part of the V18 reporting are restricted to the datasets and protocols described in Sections 6 and 7.
\end{mdframed}

\subsection{CMB compressed distance prior}
The Planck 2018 distance prior is treated as a diagnostic because it is derived under a $\Lambda$CDM modelling assumption. The MTDF evaluation against this template gives $\chi$$^{2}$/$\nu$ = 198.44 (3 d.o.f., $\chi$$^{2}$ = 595.33), reported in Table 5 below. \textbf{This large value is expected, not anomalous}: the compressed prior fixes the shape parameters \{R, $\ell$\textsubscript{A}, $\omega$\textsubscript{b}\} under a $\Lambda$CDM background expansion and growth history, and any theory that modifies either the early-universe geometry (here via the EFE injection f\textsubscript{kick}) or the late-time growth (here via $\mu$(a)) is biased away from the prior centre even when the full multipole C\textsubscript{$\ell$} spectra agree well. The honest CMB test is the full uncompressed Planck likelihood, executed with the corrected v2 implementation and reported in Section 7.3 and in companion paper \textit{MTDF\_\allowbreak{}07}~[see companion paper]: the full plik TTTEEE + lowl + lensing chain leaves the standard parameters essentially unshifted and places the theory value k\textsubscript{f} = 1 well inside the posterior (at $-$0.63$\sigma$ from the mean); the completed comparison against the matched $\Lambda$CDM control is reported in Section 7.3. The compressed-prior failure is therefore expected, structurally tied to the $\Lambda$CDM template, and superseded by the full-likelihood result. The compressed prior is excluded from the strict $\chi$$^{2}$/$\nu$ totals reported in this paper for exactly this reason.

\subsection{Gravity sector and lensing constraints (Companion Part II)}
In addition to the strict V18 cosmology likelihood set, MTDF has a dedicated gravity-sector validation programme covering
gravitational slip, Solar System screening, galaxy-galaxy lensing, strong lensing time delays, and pressure-supported systems.
These quantitative tests are reported in a separate, fully reproducible companion paper:
\textit{MTDF\_\allowbreak{}03 (Companion Part II)}~[see companion paper].
They are not yet integrated into the V18 strict dashboard totals, but they close the gravity-sector coverage gap for readers of the MTDF suite.

\begin{mdframed}[linecolor=accent,linewidth=0.5pt,backgroundcolor=callout-bg,leftline=true,rightline=true,topline=true,bottomline=true,innerleftmargin=8pt,innerrightmargin=8pt,innertopmargin=6pt,innerbottommargin=6pt]
\textbf{Gravity-sector results snapshot (outside V18 strict totals):}\begin{itemize}
  \item Gravitational slip: eta = 1 exactly.
  \item Solar System screening and PPN safety: deviations scale as O(f\_\allowbreak{}sun\^{}2), yielding very large safety margins versus current bounds.
  \item Constitutive law: uniqueness result (n = 2) and the corresponding parameter linkage.
  \item Strong lensing time delays: time-delay distances and isothermal profile behaviour consistent with H0LiCOW-style constraints.
  \item Galaxy-galaxy lensing comparisons across KiDS x GAMA and SDSS lensing samples, including a high-mass baryon completion forward prediction (Step 21).
  \item Pressure-supported systems: elliptical dispersions, satellite kinematics, and wide binaries tests (Steps 22A to 22D).
\end{itemize}
\end{mdframed}

\textbf{Scope note:} The gravity-sector paper has its own traceability chain and does not alter the strict totals in Section 7.

\section{Parameter set and provenance}
\label{sec:sec5}
The V19 workbook locks the provenance of the fundamental parameters. Table 1 lists the minimal fundamental set used throughout the strict pipeline. Where a quantity is derived, the derivation path and source DOI are recorded in the workbook.

\begin{table}[H]
\centering
\footnotesize
\begin{tabularx}{\linewidth}{l>{\RaggedRight\arraybackslash}X>{\RaggedRight\arraybackslash}X>{\RaggedRight\arraybackslash}X>{\RaggedRight\arraybackslash}X>{\RaggedRight\arraybackslash}X>{\RaggedRight\arraybackslash}X}
\toprule
\textbf{Symbol} & \textbf{Value (SI)} & \textbf{Uncertainty (SI)} & \textbf{Unit} & \textbf{Also} & \textbf{Established From} & \textbf{DOI} \tabularnewline
\midrule
$\alpha$ & 1.30 & 0.26 & dimensionless &  & Void expansion enhancement (KBC 2013 local underdensity; Hoscheit \& Barger 2018 $\Lambda$LTB model; counter: Kenworthy et al. 2019) & {\small\protect\path{10.1088/0004-637X/775/1/62}} \tabularnewline
$\beta$ & 7.00E23 & 0.09E23 & m & 22.7 Mpc & Void radius distribution (Pan et al. 2012 median \ {}24 Mpc; Hamaus et al. 2014 universal density profile) & {\small\protect\path{10.1111/j.1365-2966.2011.20197.x}} \tabularnewline
$\tau$ & 13.0 & 0.2 & Gyr &  & Age synchronisation argument: $\tau$ $\approx$ age of universe; $\gamma$ = 1/\allowbreak{}$\tau$ & {\small\protect\path{10.1051/0004-6361/201833910}} \tabularnewline
$\beta$\_\allowbreak{}eos & 0.573 & 0.012 & dimensionless &  & (Universal critical amplitudes) QCD, Landau-Ginzburg, and condensed matter & {\small\protect\path{10.1103/PhysRevD.85.054503}} \tabularnewline
\bottomrule
\end{tabularx}
\end{table}

Table 1. Fundamental MTDF parameters as locked in Params\_\allowbreak{}Fundamental (V19 workbook). The strain-lever length scale $\beta$ is also shown in megaparsecs for convenience.The V18 strict protocol distinguishes fundamental parameters from observational anchors. Anchors set scales or calibration conditions and are excluded from validation scoring.

\begin{table}[H]
\centering
\footnotesize
\begin{tabularx}{\linewidth}{l>{\RaggedRight\arraybackslash}X>{\RaggedRight\arraybackslash}X>{\RaggedRight\arraybackslash}X>{\RaggedRight\arraybackslash}X>{\RaggedRight\arraybackslash}X>{\RaggedRight\arraybackslash}X}
\toprule
\textbf{Symbol} & \textbf{Value (SI)} & \textbf{Uncertainty (SI)} & \textbf{Unit} & \textbf{Also} & \textbf{Established From} & \textbf{DOI} \tabularnewline
\midrule
$\delta$\textsubscript{bf} & 0.150 & 0.030 & 1 &  & Fixed post hoc from cluster total baryon mass fractions at M\textsubscript{500} $\sim$ 10\textsuperscript{14} M\textsubscript{$\odot$} (Giodini+2009; Gonzalez+2013); a first-principles derivation is not claimed. See \S{}5.1 note. & {\small\protect\path{10.1088/0004-637X/703/1/982;}} {\small\protect\path{10.1088/0004-637X/778/1/14}} \tabularnewline
z\textsubscript{t} & 0.74 & 0.10 & 1 &  & Farooq \& Ratra 2013 (z\_\allowbreak{}da = 0.74 $\pm$ 0.05) & {\small\protect\path{10.1088/2041-8205/766/1/L7}} \tabularnewline
E & 9.10E-10 & 0.40E-10 & Pa &  & Derived modulus: E = (2/\allowbreak{}$\alpha$$^{2}$) $\rho$\textsubscript{c} c$^{2}$ & {\small\protect\path{10.1051/0004-6361/201833910}} \tabularnewline
$\kappa$ & 0.00102 & 0.00003 & 1 &  & Theory-motivated ($\kappa$ $\approx$ f\textsubscript{kick}/3), consistent with but not measured by FIRAS bounds (see Companion Part III) & {\small\protect\path{10.1086/178173}} \tabularnewline
H\textsubscript{0} & 2.27E-18 & 4.86E-20 & s\textsuperscript{-1} & 70.0 $\pm$ 1.5 km s\textsuperscript{-1} Mpc\textsuperscript{-1} & Freedman 2021 TRGB (H$_{0}$ = 69.8 $\pm$ 0.6) & {\small\protect\path{10.3847/1538-4357/ac0e95}} \tabularnewline
\bottomrule
\end{tabularx}
\end{table}

Table 2. Observational anchors and derived quantities used in the V18 strict protocol. These values may appear in equations but are excluded from validation scoring totals. The elastic modulus E is derived from $\alpha$ and $\rho$\textsubscript{c}; the photon-stress coupling $\kappa$ is a theory-motivated, FIRAS-bounded value structurally related to the early field energy fraction f\textsubscript{kick} (see Companion Part III). Here $\rho$\textsubscript{c} denotes the background critical-density normalisation used by the gravity-sector derivation; this is distinct from the workbook local-expansion anchor H\textsubscript{0} used in validation bookkeeping.\begin{mdframed}[linecolor=accent,linewidth=0.5pt,backgroundcolor=callout-bg,leftline=true,rightline=true,topline=true,bottomline=true,innerleftmargin=8pt,innerrightmargin=8pt,innertopmargin=6pt,innerbottommargin=6pt]
\subsubsection{H\textsubscript{0} convention and predictions}
MTDF predicts a global expansion rate H\textsubscript{0,global} = 67.4 $\pm$ 0.5 km/\allowbreak{}s/\allowbreak{}Mpc (Planck-consistent). A locally inferred value near 73 km/\allowbreak{}s/\allowbreak{}Mpc is consistent with MTDF through the conditional local-environment route described in Section 3.3, contingent on a sufficiently deep local underdensity, rather than through a fixed multiplier. The relation S\textsubscript{0} = $\alpha$$\sqrt{\,}$$\Omega$\textsubscript{DE} = 1.30 $\times$ $\sqrt{\,}$0.6953 = 1.084 is an algebraic identity following from the definition of the elastic modulus E; it is noted here for its numerical proximity to the local-to-global expansion ratio, not as the physical mechanism producing the local value. The value H\textsubscript{0} = 70 $\pm$ 1.5 km/\allowbreak{}s/\allowbreak{}Mpc used as a normalization in the V19 workbook is a legacy averaging convention anchored to Freedman 2021 TRGB (69.8 $\pm$ 0.6) and does not represent the MTDF prediction. The elastic modulus E = (2/\allowbreak{}$\alpha$$^{2}$)$\rho$\textsubscript{c}c$^{2}$ in Table 2 is evaluated at the Planck-calibrated $\rho$\textsubscript{c} (H\textsubscript{0} = 67.4), not at the workbook convention of 70.
\end{mdframed}

No MTDF parameter value, frozen parameter set, validation threshold, or reported test result in this paper departs from the values locked in the V19 workbook.

\subsection{Provenance classification}
The four fundamental parameters do not share a single epistemic status, and we classify them explicitly rather than describing the set with a single label. The strain-lever length scale $\beta$ is anchored to a statistical property of an independent dataset, the distribution of void radii. The coupling $\alpha$ is treated as a frozen phenomenological coupling: it is motivated by void-dynamics arguments and is consistent with the local-underdensity literature, but it is not independently measured in the present version. The void-flow inputs that motivated its value combine measurements of different volumes and methods, and applied self-consistently to a single structure those data return a substantially smaller value; its strongest independent support is the galaxy-sector inversion of Section 4, which is itself sensitive to the stellar mass-to-light treatment. We therefore do not present $\alpha$ as an independently measured anchor, and the contested interpretation of the local underdensity (Kenworthy et al. 2019, cited as a live objection) bears directly on it. The relaxation timescale $\tau$ is fixed by an age-synchronisation ansatz, $\tau$ $\approx$ t$_{0}$; this introduces a coincidence question structurally analogous to the why-now problem of the cosmological constant, and we treat $\tau$ as a falsifiable assumption rather than a derivation. The equation-of-state parameter $\beta$\_\allowbreak{}eos is a phenomenological fit, motivated by but not derived from universal critical amplitudes in QCD and condensed-matter systems.

The accurate economy claim for MTDF is therefore: one parameter anchored to an independent observation ($\beta$), one frozen coupling motivated by void dynamics but not independently measured ($\alpha$), one synchronised to the cosmic age ($\tau$), and one fitted ($\beta$\_\allowbreak{}eos), with zero per-observable tuning thereafter. This formulation is used consistently in this version.

\subsection{Notes on interpretation}
\begin{itemize}
  \item \textbf{$\alpha$} controls the coupling between baryonic matter and the effective field stress contribution.
  \item \textbf{$\beta$} is the fundamental strain-lever length scale (\(\tilde{S} = \beta\,\nabla\xi\)): it converts displacement gradients to strain and governs transition behaviour across scales. The former ``coherence/\allowbreak{}correlation length'' label is retired (the frozen action carries no mediator mass 1/\allowbreak{}$\beta$ and no filter scale).
  \item \textbf{$\tau$} is a relaxation timescale that sets how quickly field configurations equilibrate.
  \item \textbf{$\beta$\textsubscript{eos}} enters the early epoch mapping used for the distance prior diagnostic.
  \item \textbf{$\delta$\textsubscript{bf}} enters the dynamic equation of state (E146) as the amplitude of an environmental baryon coupling f\textsubscript{env}(z). Its workbook value $\delta$\textsubscript{bf} = 0.150 $\pm$ 0.030 is consistent with cluster total baryon mass fractions measured at intermediate mass (M\textsubscript{500} $\sim$ 10\textsuperscript{14} M\textsubscript{$\odot$}), which fall in the range \~{}0.12--0.15 (Giodini et al. 2009, DOI {\small\protect\path{10.1088/0004-637X/703/1/982;}} Gonzalez, Sivanandam, Zabludoff \& Zaritsky 2013, DOI {\small\protect\path{10.1088/0004-637X/778/1/14).}} These values lie below the cosmic baryon fraction $\Omega$\textsubscript{b}/$\Omega$\textsubscript{m} = 0.157 inferred from Planck 2018 (DOI {\small\protect\path{10.1051/0004-6361/201833910);}} in the MTDF framework this deficit is interpreted as an environmental effect on the stress-tensor response, and $\delta$\textsubscript{bf} is treated as an observational anchor from cluster-scale baryon accounting rather than a fitted parameter. A derivation of $\delta$\textsubscript{bf} from first principles is not claimed in this paper and remains open for future work.
\end{itemize}

\subsection{Derived and auxiliary quantities}
Several quantities appear throughout analytic discussions as shorthand summaries of the fundamental parameter set. In the strict V18 artefacts these are treated as derived outputs rather than independently tuned parameters. Examples include:

\begin{itemize}
  \item \textbf{Critical acceleration} \(a_\mathrm{crit}\), used to indicate the low acceleration regime where departures from the classical limit become relevant.
  \item \textbf{Elastic modulus} \(E = (2/\alpha^2)\,\rho_c\,c^2\), an elastic modulus scale connecting energy density and stress in the field sector. Derived from $\alpha$ and the background critical density.
  \item \textbf{Background strain} \(S_0 = \sqrt{2\,\Omega_\mathrm{DE}\,\rho_c\,c^2/E}\). Substituting the expression for \(E\), the background strain simplifies algebraically to the exact identity \(S_0 = \alpha\sqrt{\Omega_\mathrm{DE}}\). This connects the stress-matter coupling directly to the cosmological energy budget without additional assumptions.
  \item \textbf{Relaxation rate} \(\gamma\), often approximated as \(\gamma \approx 1/\tau\) when connecting stress build up to cosmic time.
  \item \textbf{Propagation speed} \(c_s\), an effective speed for stress perturbations in reduced equations.
  \item \textbf{Effective viscosity} \(\eta\), sometimes written as \(\eta \sim E\tau\) in viscoelastic analogies, used only for intuition unless a specific observable requires it.
\end{itemize}

\subsection{Activation summary}
MTDF behaviour is governed by scale, the finite range of the stress response, and environment-dependent relaxation.

\begin{table}[H]
\centering
\footnotesize
\caption{Table 3. Regime-dependent activation of MTDF effects.}
\begin{tabularx}{\linewidth}{l>{\RaggedRight\arraybackslash}X>{\RaggedRight\arraybackslash}X}
\toprule
\textbf{Regime} & \textbf{Dominant control} & \textbf{Expected activity} \tabularnewline
\midrule
Solar System / laboratory & $\beta$-scale rigidity; individual sources do not generate a stress field & Suppressed (Newtonian) \tabularnewline
Galaxies (1--30 kpc) & Strain-lever scale ($\beta$) + low-acceleration response & Active \tabularnewline
Low-z void environments & Strain-lever scale ($\beta$) + relaxation ($\tau$) + stress depletion & Active \tabularnewline
High-z universe / early CMB & Degeneracy with $\Lambda$CDM (field not yet developed) & Effectively hidden \tabularnewline
\bottomrule
\end{tabularx}
\end{table}

\section{Validation architecture (V18 strict protocol)}
\label{sec:sec6}
The strict protocol is implemented by the V18 dashboard and workbook configuration. Pillars are treated as atomic tests. Each pillar belongs to one of four methodological roles:

\begin{itemize}
  \item \textbf{Calibration anchors (A)}: used to establish one or more fundamental parameters. Excluded from strict totals.
  \item \textbf{Benchmarks (B)}: phenomenological regularities reported as z score checks. Not used to compute strict vector totals.
  \item \textbf{Validation targets (V)}: independent observables used for strict scoring, especially the vector likelihood pillars.
  \item \textbf{Diagnostics (D)}: consistency checks, including compressed priors, used to guide development but excluded from strict totals.
\end{itemize}

\begin{mdframed}[linecolor=accent,linewidth=0.5pt,backgroundcolor=callout-bg,leftline=true,rightline=true,topline=true,bottomline=true,innerleftmargin=8pt,innerrightmargin=8pt,innertopmargin=6pt,innerbottommargin=6pt]
\textbf{Strict totals definition.} In V18, the headline ``strict $\chi$$^{2}$/$\nu$'' is computed from the combined vector validation set excluding the CMB distance prior. Scalar pillars are reported separately through z scores and scalar $\chi$$^{2}$ tallies, primarily as consistency checks.
\end{mdframed}

Model comparisons in the dashboard include several non MTDF models. For many of these, public vector based likelihood fits in exactly the same configuration are not available. The dashboard therefore tracks coverage and pass fraction separately from strict $\chi$$^{2}$/$\nu$ for those models, avoiding fabricated combined statistics.

All likelihood-based validation runs conducted under the V18 strict protocol are required to satisfy explicit numerical convergence criteria, monitored using the multivariate Gelman--Rubin statistic (R$-$1) together with per-parameter diagnostics such as effective sample size (ESS) and identification of the slowest-mixing parameters. These diagnostics are used solely to establish numerical reliability and are not included in the computation of strict validation totals.

\section{Current results (vector likelihoods and strict totals)}
\label{sec:sec7}
This section reports the current outcomes as rendered in the V75 dashboard artefacts. All values below are taken from the dashboard summary tables and are therefore the canonical reportable values for this paper.

\begin{table}[H]
\centering
\footnotesize
\begin{tabularx}{\linewidth}{l>{\RaggedRight\arraybackslash}X>{\RaggedRight\arraybackslash}X>{\RaggedRight\arraybackslash}X}
\toprule
\textbf{Likelihood set} & \textbf{$\nu$} & \textbf{MTDF $\chi$$^{2}$} & \textbf{MTDF $\chi$$^{2}$/$\nu$} \tabularnewline
\midrule
Scalar (14 counted) & 14 & 1.67 & 0.12 \tabularnewline
Vector & 1734 & 2634.12 & 1.5191 \tabularnewline
Combined & 1748 & 2635.78 & 1.5079 \tabularnewline
\textbf{Strict (excl. CMB)} & \textbf{1741} & \textbf{2036.76} & \textbf{1.1699} \tabularnewline
\bottomrule
\end{tabularx}
\end{table}

Table 4. Summary statistics from the V75 dashboard. The strict total (excluding the CMB distance prior diagnostic and the two PREDICTION-CONTESTED rows: the S$_{8}$ scalar and the DR16 growth pillar) is the headline scoring metric; the counted scalar suite is 14/14 pass at 1$\sigma$. The EFE variant differs only in the CMB diagnostic term.

\subsection{Vector pillar breakdown}
Table 5 lists the vector likelihood pillars and their $\chi$$^{2}$ contributions for MTDF. The EFE variant differs only in the CMB distance prior diagnostic term, which is shown explicitly.

\begin{table}[H]
\centering
\footnotesize
\begin{tabularx}{\linewidth}{l>{\RaggedRight\arraybackslash}X>{\RaggedRight\arraybackslash}X>{\RaggedRight\arraybackslash}X>{\RaggedRight\arraybackslash}X>{\RaggedRight\arraybackslash}X}
\toprule
\textbf{pillar\_\allowbreak{}id} & \textbf{name} & \textbf{N} & \textbf{dof} & \textbf{MTDF $\chi$$^{2}$} & \textbf{MTDF $\chi$$^{2}$/$\nu$} \tabularnewline
\midrule
P\_\allowbreak{}SNE\_\allowbreak{}PANTHEON & Pantheon+ Type Ia Supernovae & 1701 & 1700 & 2012.7 & 1.1840 \tabularnewline
P\_\allowbreak{}BAO\_\allowbreak{}DESI & DESI Y1 BAO & 12 & 12 & 15.4 & 1.2823 \tabularnewline
P\_\allowbreak{}HZ\_\allowbreak{}CC & Cosmic Chronometers H(z) & 15 & 15 & 7.0 & 0.4656 \tabularnewline
P\_\allowbreak{}GROWTH\_\allowbreak{}FSIG8 & DR16 f$\sigma$$_{8}$ Growth (PREDICTION-CONTESTED) & 4 & 4 & 3.7 & 0.9222 \tabularnewline
P\_\allowbreak{}CMB\_\allowbreak{}DIST & Planck 2018 CMB Distance Prior & 3 & 3 & 595.3 & 198.4441 \tabularnewline
\bottomrule
\end{tabularx}
\end{table}

Table 5. Vector likelihood breakdown for MTDF. The CMB entry is diagnostic and the DR16 growth entry is carried as PREDICTION-CONTESTED at the sealed amplitude; both are excluded from strict totals.

\subsection{Interpretation of the strict totals}
The strict combined reduced chi squared for MTDF is $\chi$$^{2}$/$\nu$ = 1.1699, evaluated over $\nu$ = 1741 degrees of freedom, while retaining the defining feature of MTDF: local dynamics are attributed to field stress rather than cold dark matter.

\subsection{Full-likelihood results (v2 chains, sealed posture)}

\begin{mdframed}[linecolor=accent,linewidth=0.5pt,backgroundcolor=callout-bg,leftline=true,rightline=true,topline=true,bottomline=true,innerleftmargin=8pt,innerrightmargin=8pt,innertopmargin=6pt,innerbottommargin=6pt]
\subsubsection*{Scope of the cosmological test (mandatory reading for every citation of these results)}
The v2 full-Planck run is a full Planck 2018 likelihood test of the \emph{corrected MTDF effective cosmological implementation}: the corrected covariant $\mu$(a)/\allowbreak{}EFE sector, with the high-Q carrier represented by its CDM-like limit (standard pressureless CDM as proxy). It is \emph{not} a Boltzmann implementation of all nine degrees of freedom of the frozen action. The CDM proxy is motivated by the carrier's adjudicated dust-like character; the full-action CMB perturbation implementation (unified scalar + vector perturbations) is a named future calculation (Section 2A.5). Success at this level discharges the CMB confrontation at the effective level only.
\end{mdframed}

\textbf{Run provenance.} The v2 production run (four chains, \~{}66 h) sampled the full Planck 2018 plik TTTEEE + lowl + lensing likelihood with the corrected class\_\allowbreak{}mtdf solver (covariant $\mu$ in newtonian gauge; the gauge-artifact fix is documented in the frozen source, \texttt{{\small\protect\path{perturbations.c}}}). The frozen stop rule was reached at 200,000 accepted steps per chain with R$-$1 = 0.0085 (target 0.01, satisfied); the stricter bound-level criterion R$-$1\textsubscript{CL} = 0.060 was not reached, so 95 percent bounds carry extra sampling noise while means and peaks are solid. The only free MTDF parameter is k\textsubscript{f}; $\alpha$, $\beta$\textsubscript{eos}, and $\tau$ remain frozen.

\textbf{Sealed reporting set (k\textsubscript{f} free).}
\begin{itemize}
  \item $\sigma$\textsubscript{8} = 0.8265 $\pm$ 0.0062, against the same-pipeline $\Lambda$CDM baseline 0.8101: a shift of +0.0164 (+1.91$\sigma$ in posterior units).
  \item k\textsubscript{f} = 0.685 $\pm$ 0.500: consistent with 0 and with the theory value 1 (at $-$0.63$\sigma$). Planck alone still cannot pin the EFE amplitude; consistency, not support.
  \item H\textsubscript{0} = 67.44 $\pm$ 0.55 ($\Lambda$CDM: 67.38): the CMB-side H\textsubscript{0} is unchanged.
  \item $\Omega$\textsubscript{m} = 0.3153 $\pm$ 0.0077; n\textsubscript{s} = 0.9643; log A = 3.0396 (all shifts against $\Lambda$CDM negligible).
\end{itemize}

\textbf{Theory-posture rider (k\textsubscript{f} = 1 importance reweighting).} $\sigma$\textsubscript{8} shifts by +0.0173 (+2.07$\sigma$) against the $\Lambda$CDM baseline, the same posture; f\textsubscript{kick} = 0.0033 and $\mu$(z\textsubscript{rec}) = 1.00002 (CMB-inert), as derived.

\textbf{Reading.} The corrected covariant $\mu$(a) enhances late growth by about +2 percent, and the full-likelihood chain confirms this end to end: $\sigma$\textsubscript{8} lands about 2$\sigma$ (posterior units) above the $\Lambda$CDM baseline. The corrected theory therefore mildly \emph{worsens} the comparison with published weak-lensing clustering amplitudes relative to $\Lambda$CDM; the earlier $\sigma$\textsubscript{8} = 0.790 relief remains dead as a gauge artifact. In native units the growth-sector posterior gives S\textsubscript{8} = 0.8473 $\pm$ 0.0138 (free) and S\textsubscript{8} = 0.8501 $\pm$ 0.0132 (k\textsubscript{f} = 1). Not affected by any of this: the galactic-sector validations (RAR/\allowbreak{}BTFR), c\textsubscript{T} = 1, $\eta$ = 1, the frozen 9-degree-of-freedom action, and the finding that k\textsubscript{f} is unconstrained by Planck. The registered growth-history signature (GH-1) and the preregistered forward-modelling programme (FWD-S8) define the only pre-registered paths by which this posture is tested or revised.

\textbf{Likelihood comparison against $\Lambda$CDM.} Against the matched $\Lambda$CDM control chain (identical binary, likelihood set, priors, and stop rules; three configuration deltas removing the MTDF theory keys), the bounded best fits give MTDF $-$log L = 1385.78 (at k\textsubscript{f} = 0.476) versus $\Lambda$CDM $-$log L = 1386.42: \textbf{$\Delta$$\chi$$^{2}$ (MTDF $-$ $\Lambda$CDM) = $-$1.29}. \emph{This is a maximum-likelihood comparison, not a Bayes factor or Bayesian evidence calculation.} With one additional free parameter (k\textsubscript{f}), a gain of one to two points is the no-information expectation ($\Delta$AIC = +0.71): the full Planck 2018 likelihood does not distinguish the corrected MTDF from $\Lambda$CDM. The control chain itself reproduces the reference Planck-$\Lambda$CDM posterior (per-sample S\textsubscript{8} = 0.830 $\pm$ 0.013), validating the pipeline end to end; the dataset-resolved comparison and conventions are given in Companion Part IV, section 2.4.4.

\section{Diagnostics, limitations, and road map}
\label{sec:sec8}

\subsection{What is validated in V18}
\begin{itemize}
  \item Consistency with Pantheon+ supernova distances under analytic marginalisation of the absolute magnitude nuisance parameter (as noted in Pillar\_\allowbreak{}Vector\_\allowbreak{}Config).
  \item Consistency with DESI Y1 BAO distance ratios across the configured bins.
  \item Consistency with cosmic chronometer H(z) compilation under the V18 handling of uncertainties.
  \item Consistency with DR16 f$\sigma$$_{8}$ growth points at the sealed amplitude (V19; the earlier V18-era fitted-nuisance treatment is retired).
\end{itemize}

\subsection{What is not yet validated in V18}
\begin{itemize}
  \item Full uncompressed CMB temperature and polarisation spectra likelihoods.
  \item High resolution non linear structure formation with baryonic feedback across cosmic time.
  \item Gravitational lensing power spectra and cosmic shear cross correlations treated with full systematics models. Initial galaxy-galaxy lensing and strong lens time-delay consistency checks are reported separately in \textit{MTDF\_\allowbreak{}03 (Companion Part II)}~[see companion paper].
\end{itemize}

These items are not omissions, they are explicit future tests that can refute MTDF. A high publication standard requires that the paper states these limits openly and assigns them concrete next steps.

\section{Falsifiable predictions beyond the dashboard}
\label{sec:sec9}
MTDF makes claims that can be subjected to near term tests. The following items are framed as falsifiers: if future analyses show systematic disagreement, MTDF must be revised or rejected.

\subsubsection*{The registered falsifiable set of the frozen action}
\begin{itemize}
  \item \textbf{a\textsubscript{0} magnitude (up to $\kappa$\textsubscript{a}):} the acceleration scale is anchored, a\textsubscript{MTDF} = $\kappa$\textsubscript{a}$\sqrt{\,}$(G u\textsubscript{bg}); the open O(1) constant $\kappa$\textsubscript{a} caveats the magnitude. PREDICTION-REGISTERED.
  \item \textbf{Flat a\textsubscript{0}(z):} no redshift evolution of the acceleration scale; falsifiable with high-redshift RAR/\allowbreak{}BTFR samples. PREDICTION-REGISTERED.
  \item \textbf{Cluster-scale assembly signature:} the tau-guard is load-bearing here (over-boost moderation); an active signature at cluster scale. PREDICTION-REGISTERED.
  \item \textbf{Lensing-side redshift evolution:} falsifiable with lensing surveys. PREDICTION-REGISTERED.
  \item \textbf{Isothermal-lens prediction:} the deep-MOND logarithmic potential forces time-delay lenses toward isothermal ($\gamma$' $\to$ 2) at the Einstein radius, predicting reduced $\gamma$' scatter in LSST-era samples and a small, signed few-percent low bias in lensing-inferred H\textsubscript{0} when $\Lambda$CDM composite profiles are fitted. PREDICTION-REGISTERED (dated at the H-R6 gate record).
  \item \textbf{Galaxy-scale assembly test:} CONDITIONAL, with the revival dataset named (WALLABY-era young dwarfs with resolved rotation curves and star-formation histories).
  \item \textbf{Wide-binary EFE-modulated onset:} CONTESTED (see Section 9.3, item 3).
  \item \textbf{Growth history (GH-1):} the registered f$\sigma$\textsubscript{8} trajectories; held-out growth data confirming the $\Lambda$CDM trajectory below z $\approx$ 0.74 at finer precision kills the corrected-$\mu$(a) growth branch. PREDICTION-REGISTERED.
\end{itemize}

\begin{enumerate}
  \item \textbf{CMB spectra test:} reproduce the full Planck likelihood without relying on compressed priors or $\Lambda$CDM derived distance priors.
  \item \textbf{Weak lensing consistency:} match cosmic shear and full galaxy-galaxy lensing statistics using the same parameter set, including baryon distribution assumptions. A first set of galaxy-galaxy lensing and strong lens time-delay checks is reported in \textit{MTDF\_\allowbreak{}03 (Companion Part II)}~[see companion paper], but the full cosmic shear programme remains open.
  \item \textbf{Cluster scale dynamics:} reproduce lensing mass maps and baryonic tracers in merging clusters without invoking collisionless dark matter, using transparent hydrodynamic assumptions.
  \item \textbf{Time domain probes:} test whether field relaxation implies measurable signatures in transient structure growth or environment dependent scaling relations.
  \item \textbf{Low-redshift transition-scale piecewise break:} if independent reanalyses of future SNe samples (DES-SN5YR, Rubin LSST Year 1) do not recover a piecewise transition in the range z = 0.035--0.045, the finite-range interpretation of the $\beta$ = 22.7 Mpc parameter is falsified.
\end{enumerate}

\subsection{Supplementary quantitative falsification thresholds}
The table below summarises additional quantitative thresholds that were proposed in earlier exploratory work for specific redshift and environment dependent effects. These items are not part of the strict V18 totals, but they remain useful as explicit pass fail criteria for targeted future analyses.

\begin{table}[H]
\centering
\footnotesize
\begin{tabularx}{\linewidth}{l>{\RaggedRight\arraybackslash}X>{\RaggedRight\arraybackslash}X}
\toprule
\textbf{Test} & \textbf{Expectation} & \textbf{Criterion} \tabularnewline
Environment-dependent $\Delta$z & Correlation coefficient r $\approx$ 0.34 & r < 0.20 ($\pm$3$\sigma$ of $\kappa$ uncertainty) would falsify \tabularnewline
Wavelength dependence & $\delta$z\textsubscript{X/\allowbreak{}opt} $\approx$ 0.02 at z $\approx$ 7 & $\delta$z < 0.005 would falsify \tabularnewline
Isotropy & Dipole < 5 percent & Dipole > 10 percent would falsify \tabularnewline
Time evolution & Strain growth with structure formation & No trend would falsify \tabularnewline
\bottomrule
\end{tabularx}
\end{table}

\subsection{Indicative observation timeline}
This list is an indicative sequence of observational tests that could help separate genuine field effects from survey systematics. The exact dates will depend on mission schedules and analysis readiness.

\begin{itemize}
  \item 2024--2026 -- DESI spectroscopic redshift clustering and growth analyses (ongoing)
  \item 2026 - Euclid weak-lensing correlation
  \item 2027 - Roman Space Telescope mapping
  \item 2028 - Multi-wavelength differential tests
  \item 2029 - Strain-evolution measurement
  \item 2030 - Distance-scale recalibration
\end{itemize}

\subsection{Hard falsification conditions}
The following three conditions, if established by future data, would each independently falsify a core prediction of MTDF. They are stated as explicit kill conditions rather than as tension thresholds.

\begin{enumerate}
  \item \textbf{No environment-dependent supernova signal.} If a controlled sample of approximately 3,000 or more Type Ia supernovae, matched in stretch, colour, and host-mass, shows no statistically significant (>3$\sigma$) correlation between distance residuals and void boundaries at z < 0.05, the local stress-depletion mechanism is ruled out.
  \item \textbf{No sharp z $\approx$ 0.04 transition.} If the void--supernova cross-correlation signal does not exhibit a sharp cutoff near z $\approx$ 0.04, matching the measured extent of the local underdensity (the conjectured link to $\beta$ = 22.7 Mpc is discussed in Section 10.1), the finite-range prediction is falsified. A gradual fade or no transition would be inconsistent with the framework.
  \item \textbf{Wide-binary onset (contested data, stated prediction).} The frozen action predicts a specific external-field-modulated deviation in wide binaries: the deep-MOND onset for a 1.5 M\textsubscript{$\odot$} system sits near 8,600 AU, but because the Galactic field a\textsubscript{gal} $\approx$ a\textsubscript{MTDF} places such binaries in the external-field quasi-Newtonian regime, MTDF predicts an EFE-suppressed deviation, neither the full deep-MOND boost nor strictly Newtonian behaviour. The current data are genuinely contested (Chae 2023/\allowbreak{}24 versus Banik et al. 2024 / Pittordis \& Sutherland 2023), so this row is a live falsifiable test in both directions: a confirmed full-boost signal or a confirmed exact-Newtonian null at high precision would each falsify the stated prediction.
\end{enumerate}

These falsifiers are independent and non-degenerate: satisfying one does not guarantee the others.

\section{Discriminating observational signatures}
\label{sec:sec10}
A central motivation of the Mesche Tensor Dynamics Framework is not only its ability to reproduce current precision cosmological constraints, but its capacity to generate testable discriminators that differ structurally from those implied by the baseline $\Lambda$CDM framework. The following signatures are introduced as falsifiable observational tests that arise naturally from the MTDF field dynamics and finite-range structure.

\subsection{Environment-dependent supernova residuals with redshift localisation}
MTDF implies that standardised Type Ia supernova distance residuals may exhibit a weak but systematic dependence on large-scale environment, such as proximity to cosmic voids or low-density regions. This effect arises from the finite-range response of the underlying stress field and is therefore expected to be concentrated at very low redshift, with a rapid suppression beyond a characteristic scale. Crucially, the framework does not predict a uniform global environment coefficient across the full Hubble diagram; rather, it predicts a localised transition confined to the low-redshift finite-range regime.

In particular, the framework predicts a localised environment-linked residual modulation at low redshift, with a transition near z $\approx$ 0.04 (comoving distance \~{}170 Mpc). This scale is not the field's strain-lever length scale $\beta$ $\approx$ 23 Mpc, which is roughly seven times smaller; rather, it matches the measured radial extent of the local underdensity in cluster-based reconstructions (B\"ohringer et al. 2020). A merged analysis of 4,188 supernovae (Pantheon+, ZTF DR2, Foundation DR1) is consistent with this structure: the constant-$\gamma$ global model is null ($\Delta$$\chi$\textsuperscript{2} < 1.4), while a piecewise model with a transition at z $\approx$ 0.04 yields $\Delta$$\chi$\textsuperscript{2} = 36--39 across three independent void catalogues (nominal p < 0.001, diagonal-covariance merged sample), with SALT2 population controls shifting the result by less than 0.3$\sigma$. Beyond the transition the fitted high-redshift coefficient is small and consistent with zero (the merged-sample value is approximately -0.003 at the \~{}1$\sigma$ level), as expected if the environment signal is confined to the low-redshift finite-range regime; the break is localised less sharply in the merged diagonal-covariance sample than at z = 0.04 exactly, and its precise scale is left to the full-covariance analysis. This specific redshift localisation constitutes a discriminating observational signature, since the baseline $\Lambda$CDM framework does not introduce an environment-coupled mechanism at fixed supernova standardisation, nor does it predict a preferred redshift cutoff for residual correlations. The presence of such a catalogue-independent low-redshift transition near the predicted scale provides a direct falsification channel. Full details are reported in the environmental companion (MTDF\_\allowbreak{}02, \S{}2.3).

\subsection{Interpretation of the Hubble tension as an environmental stratification}
Within MTDF, the commonly discussed Hubble tension is interpreted not primarily as a conflict between early- and late-universe physics, but as a manifestation of environmental stratification between globally averaged expansion estimates and locally weighted measurements performed in under-dense regions.

This interpretation leads to a falsifiable prediction: independent determinations of the Hubble parameter, when conditioned on large-scale environment, should exhibit a systematic separation consistent with the framework's stress-field response. Importantly, this prediction concerns relative environmental behaviour, rather than the absolute numerical value of H\textsubscript{0}, and can therefore be tested without invoking additional early-universe modifications.

\subsection{Linked direction of growth and early-structure signatures}
The corrected covariant growth sector fixes the direction of the MTDF growth signature: late-time linear growth is enhanced, by approximately +2.5 percent at z = 0 at fixed parameters, with the excess already at the \~{}2 percent level at z = 1.5 and becoming more pronounced below the $\mu$(a) transition near z\textsubscript{t} = 0.74 (the registered GH-1 trajectories freeze this prediction). An earlier version of this section predicted a suppression of late-time clustering; that direction was a gauge artifact of the v1 implementation and is formally retired. The enhanced-growth direction is the falsifiable statement: held-out growth data confirming the $\Lambda$CDM trajectory below z $\approx$ 0.74 at finer precision kills the corrected-$\mu$(a) growth branch.

A second, distinct directional signature concerns early structure. The non-analytic low-Q response of the declared A0 cubic (the same declared postulate that yields the RAR and the baryonic Tully-Fisher relation, with that ancestry and not a fitted early-universe knob) permits earlier nonlinear collapse than the standard picture, which points in the direction of the massive high-redshift galaxy candidates reported by JWST. These are two separate mechanisms and this suite never conflates them: the late linear enhancement from $\mu$(a) is active below z $\approx$ 0.74 and cannot produce structure at z \~{} 10, while any cluster-scale claim built on the early nonlinear response carries the cluster-only tau-guard explicitly named as its guard. The existence or absence of these two directional signatures, each with its own dataset, provides the discriminating avenue.

\subsection{Environment dependence of residual galaxy-scale dynamics}
At the level of individual galaxies, MTDF further implies that residuals at fixed baryonic structure may correlate with large-scale environment, such as void depth or filament membership. This expectation follows from the finite-range response of the stress field and differs qualitatively from both halo-based $\Lambda$CDM modelling, in which galaxy dynamics are primarily local, and acceleration-based modifications, in which residuals depend only on local kinematic thresholds.

Such correlations, if present, would appear as secondary trends when galaxy samples are conditioned on environment. Their absence would provide a clear falsification channel for this aspect of the framework.

\subsection{Scope and intent}
These signatures are presented as a programme of discriminating tests, not as claims of confirmation. Current datasets may show suggestive trends in some of these directions, but no statistical significance is asserted here. The primary purpose of this section is to define observational behaviours that follow directly from the MTDF framework and that can be tested independently of parameter fitting or numerical optimisation.

\section{Reproducibility and artefact traceability}
\label{sec:sec11}
The V18 era evaluation is designed to be auditable. The following items define a reproducible path from the artefacts to the reported tables:

\begin{enumerate}
  \item The workbook is the canonical source for parameter values, units, provenance strings, and vector pillar configuration.
  \item The dashboard HTML is the canonical rendering of the strict totals and pillar by pillar contributions.
  \item The validation scripts load the workbook, load vector datasets by file path, compute residuals, and assemble $\chi$$^{2}$ statistics under strict inclusion rules.
  \item The gravity-sector companion paper \textit{MTDF\_\allowbreak{}03 (Companion Part II)}~[see companion paper] provides reproducible gravity and lensing tests not yet integrated into the dashboard strict totals, with its own fixed datasets and scripts.
  \item Validation includes an independent GPU re-implementation (\textit{MTDF\_\allowbreak{}07}~[see companion paper]) which surfaced and documented a factor-2.49 coding error in an intermediate Convention-C step. The error has been corrected in all reported results; the disclosure is retained in MTDF\_\allowbreak{}07 as a reproducibility record.
\end{enumerate}

\begin{mdframed}[linecolor=accent,linewidth=0.5pt,backgroundcolor=callout-bg,leftline=true,rightline=true,topline=true,bottomline=true,innerleftmargin=8pt,innerrightmargin=8pt,innertopmargin=6pt,innerbottommargin=6pt]
\textbf{Traceability rule.} Any number reported as part of the V18 strict protocol in this master paper must be traceable to a cell in DB\_\allowbreak{}Workbook\_\allowbreak{}STRICT\_\allowbreak{}V19.xlsx or a value shown in Validation\_\allowbreak{}Dashboard\_\allowbreak{}V75.html. Results reported in separate companion analyses (for example the gravity-sector paper) have their own traceability chain and are cited here for context only.
\end{mdframed}

\section{References (DOIs)}
\label{sec:sec12}
This list is primarily assembled from the V19 workbook sources. Where a DOI represents a dataset or a calibration source, it should be cited when reproducing the corresponding pillar or parameter provenance.

{\scriptsize
\begin{thebibliography}{15}
\bibitem{ref1} \texttt{{\small\protect\path{10.1051/0004-6361/201833910}}} Field Relaxation Time ($\tau$)
\bibitem{ref2} \texttt{{\small\protect\path{10.1088/1475-7516/2025/02/021}}} P\_\allowbreak{}BAO\_\allowbreak{}DESI DESI 2024 VI
\bibitem{ref3} \texttt{{\small\protect\path{10.1088/0004-637X/775/1/62}}} Stress Tensor-Matter Coupling ($\alpha$) - KBC 2013 local underdensity
\bibitem{ref4} \texttt{{\small\protect\path{10.1111/j.1365-2966.2011.20197.x}}} Strain-lever length scale ($\beta$) - Pan et al. 2012 void radius distribution
\bibitem{ref5} \texttt{{\small\protect\path{10.1093/mnras/stt1082}}} P\_\allowbreak{}HZ\_\allowbreak{}CC Melia \& Maier 2013
\bibitem{ref6} \texttt{{\small\protect\path{10.1103/PhysRevD.103.083533}}} P\_\allowbreak{}GROWTH\_\allowbreak{}FSIG8 eBOSS 2021
\bibitem{ref7} \texttt{{\small\protect\path{10.1103/PhysRevD.106.043526}}} Smith, Lucca, Poulin \& Franco Abell\'an 2022, ``Hints of Early Dark Energy in Planck, SPT, and ACT Data'' (EDE comparator benchmark)
\bibitem{ref8} \texttt{{\small\protect\path{10.1103/PhysRevD.85.054503}}} HotQCD collaboration 2012, ``Chiral and deconfinement aspects of the QCD transition'' (cited as a physics analogue for the EoS Transition parameter $\beta$\textsubscript{eos}; the exact value $\beta$\textsubscript{eos} = 0.573 is an MTDF fit, not a direct observable of this paper)
\bibitem{ref9} \texttt{{\small\protect\path{10.1103/PhysRevLett.127.161302}}} P\_\allowbreak{}CMB\_\allowbreak{}DIST Skordis \& Z\l{}o\'snik 2021
\bibitem{ref10} \texttt{{\small\protect\path{10.3847/1538-4357/ac8e04}}} P\_\allowbreak{}SNE\_\allowbreak{}PANTHEON Brout+ 2022
\bibitem{ref11} \texttt{{\small\protect\path{10.3847/2041-8213/836/2/L24}}} MOND acceleration scale a$_{0}$ (McGaugh+2016 discovery)
\bibitem{ref12} \texttt{{\small\protect\path{10.3847/1538-4357/836/2/152}}} MOND acceleration scale a$_{0}$ (Lelli+2017, 240-galaxy confirmation)
\bibitem{ref13} \texttt{{\small\protect\path{10.1088/2041-8205/766/1/L7}}} Transition redshift (z\_\allowbreak{}t) - Farooq \& Ratra 2013
\bibitem{ref14} \texttt{{\small\protect\path{10.3847/1538-4357/ac0e95}}} H$_{0}$ anchor - Freedman 2021 TRGB
\bibitem{ref15} \texttt{{\small\protect\path{10.1051/0004-6361/201936400}}} Local underdensity reconstruction - B\"ohringer, Chon \& Collins 2020
\end{thebibliography}
}

\section{Data and code availability}
The complete MTDF reproducibility package, including the modified CLASS solver
(\texttt{class\_\allowbreak{}mtdf}), the V19 strict validation workbook, the
\texttt{{\small\protect\path{run_validate.py}}} pipeline, the gravity-sector artefact manifest,
and the Phase 2 CosmoPower emulator cache, is archived on Zenodo
(concept DOI \href{https://doi.org/10.5281/zenodo.19741058}{{\small\protect\path{10.5281/zenodo.19741058}}};
this release, \texttt{v1.1.7}, DOI
\href{https://doi.org/10.5281/zenodo.21525855}{{\small\protect\path{10.5281/zenodo.21525855}}})
and mirrored on GitHub at
\href{https://github.com/IMesche/MTDF}{github.com/\allowbreak{}IMesche/\allowbreak{}MTDF}.
The \texttt{v1.1.7} Git tag corresponds to the archived Zenodo release. All
numerical results reported here can be regenerated from the shipped workbook
and scripts following the procedure in the top-level \texttt{{\small\protect\path{README.md}}}.


\end{document}
